% Plantilla de póster · Statistical DataPulse 2026
% Facultad de Estadística, Universidad Santo Tomás · Bogotá, 11 y 12 de noviembre de 2026
%
% Tamaño final: 90 × 120 cm (ancho × alto), en vertical, como dice la
% convocatoria. Compila con pdfLaTeX (también en Overleaf).
%
% Qué editar:
%   1. Los datos del póster, en la sección «Datos del póster».
%   2. El logo de tu institución: reemplaza Logo.png o cambia \logoinstitucion.
%   3. El contenido de los bloques, entre \begin{document} y \end{document}.
%
% Gráficas: están hechas en TikZ/pgfplots (carpeta figs_tikz/, generadas por
% gen_tikz.py con los datos del estudio), así usan la misma fuente del póster.

\documentclass[final]{beamer}

% ====================
% Paquetes
% ====================

\usepackage[T1]{fontenc}
\usepackage[spanish,es-noshorthands]{babel}
\usepackage[orientation=portrait,size=custom,width=90,height=120,scale=1.2]{beamerposter}
\usepackage[sfdefault]{plex-sans}
\usepackage{plex-mono}
\usepackage{ragged2e}
\usepackage{graphicx}
\usepackage{booktabs}
\usepackage{tikz}
\usepackage[most]{tcolorbox}
\usepackage{anyfontsize}
\usepackage{pgfplots}
\addto\captionsspanish{\renewcommand{\tablename}{Tabla}}
\pgfplotsset{compat=1.18}

% ====================
% Datos del póster
% ====================

% Un título de hasta tres líneas cabe; el encabezado crece y las columnas se
% ajustan solas. Usa ~ para que algo no se parta entre líneas: «D.~C.».
\title{Cómo contrata el Estado los servicios funerarios: muestreo y clustering en SECOP, 2020--2025}

% Separa los autores con comas; \inst{n} es el número de su afiliación.
\author{Edward Mora\inst{1}, Natalia Carrero\inst{1}, Harold Funeme\inst{1}}

\institute{\inst{1}\,Universidad Santo Tomás, Pregrado en Estadística}

% Logo de la institución de los autores (va en un recuadro blanco).
\newcommand{\logoinstitucion}{Logo.png}

% Correos del equipo.
\newcommand{\correocontacto}{edwardmora@usantotomas.edu.co\\nataliacarrero@usantotomas.edu.co\\haroldfuneme@usantotomas.edu.co}

% ====================
% Colores del simposio (los mismos del sitio)
% ====================

\definecolor{noche}{HTML}{000D1C}
\definecolor{naranja}{HTML}{F29200}
\definecolor{naranjatexto}{HTML}{A65400}
\definecolor{crema}{HTML}{FFE6BE}
\definecolor{cremafondo}{HTML}{FFF0DA}
\definecolor{tinta}{HTML}{111820}
\definecolor{tintamedia}{HTML}{2C3743}
\definecolor{sobrenoche}{HTML}{C3CCD6}
\definecolor{lineafuerte}{HTML}{D8D0C4}
\definecolor{gris}{HTML}{8A94A3}
% Grupos del clustering (G1 a G8)
\definecolor{g1}{HTML}{0047BB}\definecolor{g2}{HTML}{5B2A86}\definecolor{g3}{HTML}{C2410C}\definecolor{g4}{HTML}{0F766E}
\definecolor{g5}{HTML}{A16207}\definecolor{g6}{HTML}{BE185D}\definecolor{g7}{HTML}{1B7A3E}\definecolor{g8}{HTML}{0A1F44}
% Fuente de cada conclusión.
\newcommand{\fuente}[1]{\par{\footnotesize\color{tintamedia}Fuente: #1\par}}
% Tarjeta de indicador: cifra grande y explicación corta.
\newcommand{\kpi}[2]{%
  \begin{tcolorbox}[enhanced,colback=cremafondo,colframe=naranja,boxrule=0pt,toprule=0.25cm,arc=0.3cm,
    left=0.7cm,right=0.7cm,top=0.45cm,bottom=0.45cm,equal height group=kpi]
  {\fontsize{58}{62}\selectfont\bfseries\color{tinta}#1\par}\vskip0.3em
  {\footnotesize\raggedright\color{tintamedia}#2\par}
  \end{tcolorbox}}

% ====================
% Estilo
% ====================

\usefonttheme{professionalfonts}
\setbeamercolor{background canvas}{bg=white}
\setbeamercolor{normal text}{fg=tintamedia,bg=white}
\setbeamercolor{structure}{fg=tinta}
\setbeamercolor{headline}{fg=white,bg=noche}
\setbeamercolor{footline}{fg=white,bg=noche}
\setbeamercolor{caption name}{fg=naranjatexto}
\setbeamerfont{title}{size=\Huge,series=\bfseries}
\setbeamerfont{author}{size=\LARGE}
\setbeamerfont{institute}{size=\large}
\setbeamerfont{block title}{size=\huge,series=\bfseries}
\setbeamerfont{caption}{size=\small}
\setbeamerfont{caption name}{series=\bfseries}
\setbeamertemplate{navigation symbols}{}
\setbeamertemplate{caption}[numbered]
\setbeamertemplate{caption label separator}[period]

% Viñetas: cuadros naranja. Pasos numerados: círculos naranja, como en el sitio.
\setbeamertemplate{itemize item}{\textcolor{naranja}{\rule[0.18em]{0.42em}{0.42em}}}
\setbeamertemplate{itemize subitem}{\textcolor{naranja}{\rule[0.22em]{0.3em}{0.3em}}}
\setbeamertemplate{enumerate item}{%
  \tikz[baseline=(n.base)]\node[circle,fill=naranja,text=tinta,inner sep=0pt,
    minimum size=1.35em,font=\small\bfseries](n){\insertenumlabel};}

% Regla bajo el título de cada bloque: tramo naranja y línea gris.
\newcommand{\reglabloque}{%
  {\color{naranja}\rule{3.2cm}{0.24cm}}%
  {\color{lineafuerte}\rule[0.09cm]{\dimexpr\linewidth-3.2cm\relax}{0.06cm}}}

\setbeamertemplate{block begin}{%
  \par\vskip0.85cm
  {\usebeamerfont{block title}\color{tinta}\insertblocktitle\par}%
  \vskip0.45cm\reglabloque\par\vskip0.8cm
  \raggedright\setlength{\parskip}{0.6em}%
}
\setbeamertemplate{block end}{\par}

% Bloque destacado (alertblock): fondo crema con borde naranja a la izquierda.
\setbeamertemplate{block alerted begin}{%
  \par\vskip0.85cm
  \begin{tcolorbox}[enhanced,colback=cremafondo,colframe=naranja,boxrule=0pt,
    leftrule=0.3cm,arc=0.4cm,left=1.2cm,right=1.2cm,top=1cm,bottom=1cm]
  {\usebeamerfont{block title}\color{tinta}\insertblocktitle\par}%
  \vskip0.8cm
  \raggedright\setlength{\parskip}{0.6em}%
}
\setbeamertemplate{block alerted end}{\end{tcolorbox}}

% Subtítulo dentro de un bloque.
\newcommand{\subtitulo}[1]{\par\vskip0.5em{\large\bfseries\color{tinta}#1\par}\vskip0.2em}

% ====================
% Encabezado y pie
% ====================

\newlength{\sepwidth}
\newlength{\colwidth}
\setlength{\sepwidth}{0.03\paperwidth}
\setlength{\colwidth}{0.455\paperwidth}   % 3 × 0,03 + 2 × 0,455 = 1
\newcommand{\separatorcolumn}{\begin{column}{\sepwidth}\end{column}}
% Alto de las columnas: se calcula al empezar el póster con lo que dejan
% libre el encabezado (crece si el título ocupa más líneas) y el pie. Los
% \vfill entre bloques reparten el espacio sobrante para que las dos
% columnas terminen a la misma altura.
\newlength{\altocolumnas}
\newlength{\anchologo}
\setlength{\anchologo}{8cm}

\setbeamertemplate{headline}{%
  \begin{beamercolorbox}[wd=\paperwidth]{headline}
    \vskip2.4cm
    \hspace*{\sepwidth}%
    \begin{minipage}[c]{\dimexpr\anchologo+1.6cm\relax}
      \begin{tikzpicture}
        \node[fill=white,rounded corners=0.5cm,inner sep=0.8cm]
          {\includegraphics[width=\anchologo,height=\anchologo,keepaspectratio]{\logoinstitucion}};
      \end{tikzpicture}
    \end{minipage}%
    \hspace{2.2cm}%
    \begin{minipage}[c]{\dimexpr\paperwidth-2\sepwidth-\anchologo-3.8cm\relax}
      \raggedright
      {\usebeamerfont{title}\color{white}\inserttitle\par}
      \vskip1.3cm
      {\usebeamerfont{author}\color{crema}\insertauthor\par}
      \vskip0.6cm
      {\usebeamerfont{institute}\color{sobrenoche}\insertinstitute\par}
    \end{minipage}%
    \vskip2.4cm
    \noindent{\color{naranja}\rule{\paperwidth}{0.5cm}}%
  \end{beamercolorbox}%
}

\setbeamertemplate{footline}{%
  \begin{beamercolorbox}[wd=\paperwidth]{footline}
    \vskip1.6cm
    \hspace*{\sepwidth}%
    \begin{minipage}[c]{13cm}
      \includegraphics[width=\linewidth]{logo_evento_blanco.png}
    \end{minipage}%
    \hspace{2.2cm}%
    \begin{minipage}[c]{40cm}
      \raggedright\normalsize
      {\color{white}\bfseries Simposio de la Facultad de Estadística · Universidad Santo Tomás\par}
      \vskip0.25em
      {\color{sobrenoche}Bogotá, 11 y 12 de noviembre de 2026\par}
    \end{minipage}%
    \hfill
    \begin{minipage}[c]{20cm}
      \raggedleft\normalsize
      {\color{sobrenoche}\small\ttfamily CONTACTO\par}
      \vskip0.2em
      {\color{white}\small\correocontacto\par}
    \end{minipage}%
    \hspace*{\sepwidth}
    \vskip1.6cm
  \end{beamercolorbox}%
}

% ====================
% Contenido
% ====================

\begin{document}

\begin{frame}[t]
\setlength{\altocolumnas}{\dimexpr\textheight-5.1cm\relax}
\vskip0.6cm
\begin{columns}[t]
\separatorcolumn

\begin{column}{\colwidth}
\begin{minipage}[t][\altocolumnas][s]{\linewidth}

  \begin{block}{Contexto}

    Alcaldías, gobernaciones, hospitales públicos y empresas de servicios públicos contratan con funerarias, cementerios y entidades sin ánimo de lucro el servicio funerario de población vulnerable, víctimas y cuerpos no identificados. En SECOP \textbf{no hay un código que identifique estos contratos}: el UNSPSC se diligencia mal y buscar por palabras trae sobre todo ruido.

    \textbf{Por qué medirlo.} Es un gasto público anual que hoy no se puede filtrar ni vigilar en SECOP. La contratación pública ya se estudia con microdatos (Fazekas y Kocsis, 2020), pero no existía una medición del mercado funerario público que corrigiera el error de clasificación.

  \end{block}
  \vfill

  \begin{block}{Objetivos}
    \textbf{\color{tinta}Objetivo general:} caracterizar cómo contrata el Estado colombiano los servicios funerarios en SECOP I y II, 2020–2025.

    {\bfseries\color{tinta}Objetivos específicos\par}
    \begin{itemize}
      \setlength{\itemsep}{0.4em}
      \item \textbf{Identificar} los contratos funerarios entre 39.870 candidatos.
      \item \textbf{Estimar} el valor del mercado y su incertidumbre.
      \item \textbf{Modelar} qué se asocia al valor de un contrato y en qué meses se contrata.
      \item \textbf{Agrupar} las formas de contratar y comprobarlas con valor, territorio y proveedores.
    \end{itemize}
  \end{block}
  \vfill

  \begin{alertblock}{Metodología}
    \begin{enumerate}
      \setlength{\itemsep}{0.4em}
      \item \textbf{ETL:} API de SECOP I y II; 39.870 candidatos por palabra clave o UNSPSC, deduplicados y en pesos de 2025.
      \item \textbf{Muestreo estratificado:} 400 contratos en 9 estratos, pesos $w_h = N_h/n_h$, Horvitz-Thompson ($\hat{p} = \sum w_k y_k / \sum w_k$) e intervalos de Wilson.
      \item \textbf{Medición:} etiquetado doble ($\kappa$ de Cohen) y clasificador TF-IDF con regresión logística; otro modelo lo reemplaza solo si su IC del F1 queda entero por encima.
      \item \textbf{Estimador de diferencia} del total, $\hat{T} = \sum_{U} \hat{y}_k + \sum_h \frac{N_h}{n_h} \sum_{s_h} (y_k - \hat{y}_k)$, con bootstrap estratificado.
      \item \textbf{Modelos} del log del valor (lineal, mixto, gamma, cuantiles) y del conteo mensual (binomial negativa). Normalidad de residuos: $|\text{asimetría}| \leq 1$ y $\text{curtosis} \leq 3$.
      \item \textbf{ACM y HCPC}; estable si su Jaccard medio es $\geq 0{,}75$.
    \end{enumerate}
  \end{alertblock}
  \vfill

  \begin{block}{Muestreo y medición}
    \begin{columns}[t]
      \begin{column}{0.32\linewidth}\kpi{0,683}{$\kappa$ de Cohen entre revisores: sustancial, bajo el umbral fijado de 0,70.}\end{column}
      \begin{column}{0.32\linewidth}\kpi{0,956}{F1 del clasificador en validación externa (IC 95\,\%: 0,908–0,996).}\end{column}
      \begin{column}{0.32\linewidth}\kpi{4,70}{Efecto de diseño de Kish: los 400 contratos rinden como 85.}\end{column}
    \end{columns}
    \vskip0.4em
    \begin{figure}
      \centering
      % Generado por gen_tikz.py
\begin{tikzpicture}
\begin{axis}[width=0.52\linewidth,height=6.3cm,scale only axis,xmin=-1.5,xmax=101.5,ymin=-0.7,ymax=9,clip=false,
  axis y line=none,axis x line*=bottom,xtick={0,25,50,75,100},xticklabels={0\,\%,25\,\%,50\,\%,75\,\%,100\,\%},
  ticklabel style={font=\footnotesize,text=tintamedia},xmajorgrids,grid style={lineafuerte,line width=1pt},axis line style={tintamedia,line width=1.2pt},
  xlabel={Proporción funeraria en la muestra (punto) e intervalo de Wilson del 95\,\%},xlabel style={font=\footnotesize,text=tinta}]
\draw[naranja,opacity=.45,line width=6pt] (axis cs:93.98,8) -- (axis cs:100,8);
\fill[naranja,draw=white,line width=1.5pt] (axis cs:100.00,8) circle[radius=5.5pt];
\node[anchor=east,font=\scriptsize\bfseries,text=tinta] at (axis cs:-3,8) {A1 · palabra y código};
\node[anchor=west,font=\scriptsize\bfseries,text=naranjatexto] at (axis cs:106,8) {100\,\%};
\node[anchor=west,font=\tiny,text=tintamedia] at (axis cs:121,8) {60 de 2.312};
\draw[naranja,opacity=.45,line width=6pt] (axis cs:34.62,7) -- (axis cs:54.73,7);
\fill[naranja,draw=white,line width=1.5pt] (axis cs:44.44,7) circle[radius=5.5pt];
\node[anchor=east,font=\scriptsize\bfseries,text=tinta] at (axis cs:-3,7) {A2 · solo palabra};
\node[anchor=west,font=\scriptsize\bfseries,text=naranjatexto] at (axis cs:106,7) {44\,\%};
\node[anchor=west,font=\tiny,text=tintamedia] at (axis cs:121,7) {90 de 1.469};
\draw[naranja,opacity=.45,line width=6pt] (axis cs:59.23,6) -- (axis cs:81.49,6);
\fill[naranja,draw=white,line width=1.5pt] (axis cs:71.67,6) circle[radius=5.5pt];
\node[anchor=east,font=\scriptsize\bfseries,text=tinta] at (axis cs:-3,6) {A3 · solo código};
\node[anchor=west,font=\scriptsize\bfseries,text=naranjatexto] at (axis cs:106,6) {72\,\%};
\node[anchor=west,font=\tiny,text=tintamedia] at (axis cs:121,6) {60 de 976};
\draw[tintamedia,opacity=.45,line width=6pt] (axis cs:0,5) -- (axis cs:14.31,5);
\fill[noche,draw=white,line width=1.5pt] (axis cs:0.00,5) circle[radius=5.5pt];
\node[anchor=east,font=\scriptsize\bfseries,text=tinta] at (axis cs:-3,5) {B · animales};
\node[anchor=west,font=\scriptsize\bfseries,text=tintamedia] at (axis cs:106,5) {0\,\%};
\node[anchor=west,font=\tiny,text=tintamedia] at (axis cs:121,5) {23 de 23};
\draw[tintamedia,opacity=.45,line width=6pt] (axis cs:0,4) -- (axis cs:11.35,4);
\fill[noche,draw=white,line width=1.5pt] (axis cs:0.00,4) circle[radius=5.5pt];
\node[anchor=east,font=\scriptsize\bfseries,text=tinta] at (axis cs:-3,4) {C · cadáver};
\node[anchor=west,font=\scriptsize\bfseries,text=tintamedia] at (axis cs:106,4) {0\,\%};
\node[anchor=west,font=\tiny,text=tintamedia] at (axis cs:121,4) {30 de 578};
\draw[tintamedia,opacity=.45,line width=6pt] (axis cs:0,3) -- (axis cs:7.87,3);
\fill[noche,draw=white,line width=1.5pt] (axis cs:0.00,3) circle[radius=5.5pt];
\node[anchor=east,font=\scriptsize\bfseries,text=tinta] at (axis cs:-3,3) {C · cementerio};
\node[anchor=west,font=\scriptsize\bfseries,text=tintamedia] at (axis cs:106,3) {0\,\%};
\node[anchor=west,font=\tiny,text=tintamedia] at (axis cs:121,3) {45 de 2.378};
\draw[tintamedia,opacity=.45,line width=6pt] (axis cs:0,2) -- (axis cs:18.43,2);
\fill[noche,draw=white,line width=1.5pt] (axis cs:0.00,2) circle[radius=5.5pt];
\node[anchor=east,font=\scriptsize\bfseries,text=tinta] at (axis cs:-3,2) {C · osario y mixtas};
\node[anchor=west,font=\scriptsize\bfseries,text=tintamedia] at (axis cs:106,2) {0\,\%};
\node[anchor=west,font=\tiny,text=tintamedia] at (axis cs:121,2) {17 de 98};
\draw[tintamedia,opacity=.45,line width=6pt] (axis cs:0,1) -- (axis cs:13.32,1);
\fill[noche,draw=white,line width=1.5pt] (axis cs:0.00,1) circle[radius=5.5pt];
\node[anchor=east,font=\scriptsize\bfseries,text=tinta] at (axis cs:-3,1) {C · necropsia};
\node[anchor=west,font=\scriptsize\bfseries,text=tintamedia] at (axis cs:106,1) {0\,\%};
\node[anchor=west,font=\tiny,text=tintamedia] at (axis cs:121,1) {25 de 560};
\draw[tintamedia,opacity=.45,line width=6pt] (axis cs:0,0) -- (axis cs:7.14,0);
\fill[noche,draw=white,line width=1.5pt] (axis cs:0.00,0) circle[radius=5.5pt];
\node[anchor=east,font=\scriptsize\bfseries,text=tinta] at (axis cs:-3,0) {D · sin señal};
\node[anchor=west,font=\scriptsize\bfseries,text=tintamedia] at (axis cs:106,0) {0\,\%};
\node[anchor=west,font=\tiny,text=tintamedia] at (axis cs:121,0) {50 de 26.396};
\node[anchor=west,font=\scriptsize\bfseries,text=tintamedia] at (axis cs:106,8.95) {$\hat p$};
\node[anchor=west,font=\scriptsize\bfseries,text=tintamedia] at (axis cs:121,8.95) {muestra de $N$};
\end{axis}
\end{tikzpicture}

      \caption{Proporción funeraria por estrato en la muestra de 400. En total, $\hat{p} = 10{,}5$\,\% (IC 95\,\%: 10,0–11,1): uno de cada diez candidatos}
    \end{figure}
  \end{block}

\end{minipage}
\end{column}

\separatorcolumn

\begin{column}{\colwidth}
\begin{minipage}[t][\altocolumnas][s]{\linewidth}

  \begin{block}{Resultados}
    \begin{columns}[t]
      \begin{column}{0.49\linewidth}
        \begin{figure}
          \centering
          % Generado por gen_tikz.py
\begin{tikzpicture}
\begin{axis}[width=0.74\linewidth,height=9cm,scale only axis,xmin=-0.5,xmax=5.3,ymin=150,ymax=330,clip=false,
  xtick={0,1,2,3,4},xticklabels={Modelo,Top 50,Muestra,Lote,Final},ytick={150,200,250,300},yticklabels={150,200,250,300},
  ticklabel style={font=\scriptsize,text=tintamedia},axis x line*=bottom,axis y line*=left,ymajorgrids,grid style={lineafuerte,line width=1pt},
  axis line style={tintamedia,line width=1.2pt},major tick style={tintamedia},
  ylabel={Miles de millones de pesos de 2025},ylabel style={font=\scriptsize,text=tinta,align=center}]
\fill[noche] (axis cs:-0.33,150) rectangle (axis cs:0.33,262.9);
\node[anchor=south,font=\footnotesize\bfseries,text=tinta] at (axis cs:0,266.9) {262,9};
\fill[g3] (axis cs:0.6699999999999999,229.39999999999998) rectangle (axis cs:1.33,262.9);
\node[anchor=south,font=\footnotesize\bfseries,text=g3] at (axis cs:1,266.9) {--33,5};
\fill[g7] (axis cs:1.67,229.39999999999998) rectangle (axis cs:2.33,253.99999999999997);
\node[anchor=south,font=\footnotesize\bfseries,text=g7] at (axis cs:2,258.0) {+24,6};
\fill[g7] (axis cs:2.67,253.99999999999997) rectangle (axis cs:3.33,255.19999999999996);
\node[anchor=south,font=\footnotesize\bfseries,text=g7] at (axis cs:3,258.09999999999997) {+0,1};
\fill[naranja] (axis cs:3.67,150) rectangle (axis cs:4.33,254.0);
\node[anchor=south,font=\footnotesize\bfseries,text=noche] at (axis cs:4,175) {254,0};
\draw[tinta,line width=2.5pt] (axis cs:4,225.4) -- (axis cs:4,304.5); \draw[tinta,line width=2.5pt] (axis cs:3.85,225.4) -- (axis cs:4.15,225.4); \draw[tinta,line width=2.5pt] (axis cs:3.85,304.5) -- (axis cs:4.15,304.5);
\node[anchor=north west,font=\scriptsize,text=tintamedia,align=left] at (axis cs:4.25,306) {IC 95\,\%\\225,4--304,5};
\end{axis}
\end{tikzpicture}

          \caption{Del modelo al estimador de diferencia}
        \end{figure}
      \end{column}
      \begin{column}{0.49\linewidth}
        \begin{figure}
          \centering
          % Generado por gen_tikz.py
\begin{tikzpicture}
\begin{axis}[width=0.78\linewidth,height=9cm,scale only axis,xmin=2019.5,xmax=2025.5,ymin=0,ymax=68,clip=false,
  xtick={2020,...,2025},xticklabels={2020,21,22,23,24,2025},ytick={0,25,50},yticklabels={0,25,50},
  ticklabel style={font=\footnotesize,text=tintamedia},axis x line*=bottom,axis y line*=left,ymajorgrids,grid style={lineafuerte,line width=1pt},
  axis line style={tintamedia,line width=1.2pt},major tick style={tintamedia},
  ylabel={Miles de millones de pesos de 2025},ylabel style={font=\scriptsize,text=tinta}]
\fill[noche] (axis cs:2019.67,0) rectangle (axis cs:2020.33,27.83); \fill[naranja] (axis cs:2019.67,27.83) rectangle (axis cs:2020.33,44.65);
\node[anchor=south,font=\scriptsize\bfseries,text=tinta] at (axis cs:2020,45.15) {44,6};
\node[font=\tiny\bfseries,text=tinta] at (axis cs:2020,36.239999999999995) {8\,\%};
\fill[noche] (axis cs:2020.67,0) rectangle (axis cs:2021.33,19.95); \fill[naranja] (axis cs:2020.67,19.95) rectangle (axis cs:2021.33,41.68);
\node[anchor=south,font=\scriptsize\bfseries,text=tinta] at (axis cs:2021,42.18) {41,7};
\node[font=\tiny\bfseries,text=tinta] at (axis cs:2021,30.814999999999998) {16\,\%};
\fill[noche] (axis cs:2021.67,0) rectangle (axis cs:2022.33,15.36); \fill[naranja] (axis cs:2021.67,15.36) rectangle (axis cs:2022.33,36.44);
\node[anchor=south,font=\scriptsize\bfseries,text=tinta] at (axis cs:2022,36.94) {36,4};
\node[font=\tiny\bfseries,text=tinta] at (axis cs:2022,25.9) {26\,\%};
\fill[noche] (axis cs:2022.67,0) rectangle (axis cs:2023.33,12.98); \fill[naranja] (axis cs:2022.67,12.98) rectangle (axis cs:2023.33,33.96);
\node[anchor=south,font=\scriptsize\bfseries,text=tinta] at (axis cs:2023,34.46) {34,0};
\node[font=\tiny\bfseries,text=tinta] at (axis cs:2023,23.47) {31\,\%};
\fill[noche] (axis cs:2023.67,0) rectangle (axis cs:2024.33,9.97); \fill[naranja] (axis cs:2023.67,9.97) rectangle (axis cs:2024.33,35.01);
\node[anchor=south,font=\scriptsize\bfseries,text=tinta] at (axis cs:2024,35.51) {35,0};
\node[font=\tiny\bfseries,text=tinta] at (axis cs:2024,22.490000000000002) {43\,\%};
\fill[noche] (axis cs:2024.67,0) rectangle (axis cs:2025.33,9.22); \fill[naranja] (axis cs:2024.67,9.22) rectangle (axis cs:2025.33,37.88);
\node[anchor=south,font=\scriptsize\bfseries,text=tinta] at (axis cs:2025,38.38) {37,9};
\node[font=\tiny\bfseries,text=tinta] at (axis cs:2025,23.55) {51\,\%};
\fill[noche] (axis cs:2019.7,62) rectangle (axis cs:2020.0,65.5); \node[anchor=west,font=\scriptsize,text=tinta] at (axis cs:2020.05,63.75) {SECOP I};
\fill[naranja] (axis cs:2021.35,62) rectangle (axis cs:2021.65,65.5); \node[anchor=west,font=\scriptsize,text=tinta] at (axis cs:2021.7,63.75) {SECOP II (porcentaje de contratos)};
\end{axis}
\end{tikzpicture}

          \caption{Valor anual por plataforma}
        \end{figure}
      \end{column}
    \end{columns}

    \begin{columns}[t]
      \begin{column}{0.50\linewidth}
        \begin{figure}
          \centering
          % Generado por gen_tikz.py
\begin{tikzpicture}
\begin{axis}[width=0.82\linewidth,height=10cm,scale only axis,xmin=-0.95,xmax=2.95,ymin=-1.85,ymax=2.3,clip=true,
  axis x line*=bottom,axis y line*=left,xtick=\empty,ytick=\empty,axis line style={tintamedia,line width=1.2pt},
  xlabel={Eje 1 (12,6\,\%): compra institucional $\rightarrow$},ylabel={Eje 2 (9,2\,\%): bienes y planes $\uparrow$},
  label style={font=\scriptsize,text=tinta}]
\draw[tintamedia,line width=1.2pt,dash pattern=on 6pt off 5pt] (axis cs:-0.95,0) -- (axis cs:2.95,0); \draw[tintamedia,line width=1.2pt,dash pattern=on 6pt off 5pt] (axis cs:0,-1.85) -- (axis cs:0,2.3);
\fill[g1,opacity=.6] (axis cs:-0.033,-0.335) circle[radius=60.2pt];
\fill[g2,opacity=.6] (axis cs:-0.521,0.397) circle[radius=45.0pt];
\fill[g1,opacity=.6] (axis cs:-0.252,-0.125) circle[radius=35.8pt];
\fill[g2,opacity=.6] (axis cs:-0.484,0.089) circle[radius=30.7pt];
\fill[g2,opacity=.6] (axis cs:-0.293,0.183) circle[radius=22.3pt];
\fill[g1,opacity=.6] (axis cs:0.234,-0.624) circle[radius=22.1pt];
\fill[g4,opacity=.6] (axis cs:-0.521,0.526) circle[radius=21.8pt];
\fill[g1,opacity=.6] (axis cs:-0.061,-0.026) circle[radius=21.6pt];
\fill[g1,opacity=.6] (axis cs:0.316,-0.239) circle[radius=19.9pt];
\fill[g5,opacity=.6] (axis cs:0.045,-0.443) circle[radius=17.0pt];
\fill[g1,opacity=.6] (axis cs:0.181,0.360) circle[radius=14.1pt];
\fill[g7,opacity=.6] (axis cs:1.149,-0.746) circle[radius=13.4pt];
\fill[g4,opacity=.6] (axis cs:-0.488,0.221) circle[radius=12.8pt];
\fill[g1,opacity=.6] (axis cs:0.197,-0.315) circle[radius=11.8pt];
\fill[g5,opacity=.6] (axis cs:-0.406,-0.019) circle[radius=11.5pt];
\fill[g3,opacity=.6] (axis cs:-0.030,-0.354) circle[radius=11.5pt];
\fill[g2,opacity=.6] (axis cs:-0.042,0.596) circle[radius=10.5pt];
\fill[g2,opacity=.6] (axis cs:-0.175,0.492) circle[radius=10.4pt];
\fill[g5,opacity=.6] (axis cs:-0.175,-0.229) circle[radius=10.4pt];
\fill[g3,opacity=.6] (axis cs:-0.256,-0.142) circle[radius=10.1pt];
\fill[g2,opacity=.6] (axis cs:-0.135,0.184) circle[radius=9.9pt];
\fill[g2,opacity=.6] (axis cs:0.061,0.280) circle[radius=9.7pt];
\fill[g3,opacity=.6] (axis cs:0.072,-0.444) circle[radius=9.6pt];
\fill[g1,opacity=.6] (axis cs:0.171,0.017) circle[radius=9.6pt];
\fill[g3,opacity=.6] (axis cs:0.323,-0.264) circle[radius=9.6pt];
\fill[g1,opacity=.6] (axis cs:0.280,0.067) circle[radius=9.6pt];
\fill[g3,opacity=.6] (axis cs:-0.506,0.367) circle[radius=9.6pt];
\fill[g6,opacity=.6] (axis cs:1.643,1.318) circle[radius=9.4pt];
\fill[g1,opacity=.6] (axis cs:0.583,-0.528) circle[radius=9.2pt];
\fill[g6,opacity=.6] (axis cs:1.127,0.289) circle[radius=9.1pt];
\fill[g3,opacity=.6] (axis cs:-0.481,0.070) circle[radius=8.5pt];
\fill[g1,opacity=.6] (axis cs:0.228,0.128) circle[radius=8.5pt];
\fill[g1,opacity=.6] (axis cs:0.546,-0.219) circle[radius=8.3pt];
\fill[g4,opacity=.6] (axis cs:-0.036,0.700) circle[radius=8.1pt];
\fill[g1,opacity=.6] (axis cs:0.531,0.462) circle[radius=8.1pt];
\fill[g6,opacity=.6] (axis cs:0.990,0.772) circle[radius=8.1pt];
\fill[g3,opacity=.6] (axis cs:0.586,-0.547) circle[radius=7.5pt];
\fill[g2,opacity=.6] (axis cs:-0.270,0.784) circle[radius=7.5pt];
\fill[g7,opacity=.6] (axis cs:1.607,-1.318) circle[radius=7.3pt];
\fill[g3,opacity=.6] (axis cs:0.124,-0.032) circle[radius=6.9pt];
\fill[g8,opacity=.6] (axis cs:1.625,0.887) circle[radius=6.9pt];
\fill[g7,opacity=.6] (axis cs:1.750,0.790) circle[radius=6.6pt];
\fill[g5,opacity=.6] (axis cs:0.007,-0.134) circle[radius=6.6pt];
\fill[g8,opacity=.6] (axis cs:1.589,1.188) circle[radius=6.4pt];
\fill[g3,opacity=.6] (axis cs:0.431,-0.342) circle[radius=6.4pt];
\fill[g8,opacity=.6] (axis cs:1.885,1.329) circle[radius=6.4pt];
\fill[g4,opacity=.6] (axis cs:-0.250,0.005) circle[radius=6.2pt];
\fill[g3,opacity=.6] (axis cs:0.339,-0.009) circle[radius=6.2pt];
\fill[g5,opacity=.6] (axis cs:-0.444,0.290) circle[radius=6.2pt];
\fill[g6,opacity=.6] (axis cs:0.662,0.647) circle[radius=5.9pt];
\fill[g8,opacity=.6] (axis cs:1.839,1.582) circle[radius=5.9pt];
\fill[g1,opacity=.6] (axis cs:0.005,-0.411) circle[radius=5.9pt];
\fill[g3,opacity=.6] (axis cs:0.205,-0.614) circle[radius=5.6pt];
\fill[g3,opacity=.6] (axis cs:-0.267,0.765) circle[radius=5.6pt];
\fill[g3,opacity=.6] (axis cs:-0.146,-0.234) circle[radius=5.6pt];
\fill[g5,opacity=.6] (axis cs:0.504,-0.126) circle[radius=5.6pt];
\fill[g3,opacity=.6] (axis cs:-0.005,-0.131) circle[radius=5.6pt];
\fill[g3,opacity=.6] (axis cs:-0.283,0.163) circle[radius=5.3pt];
\fill[g6,opacity=.6] (axis cs:0.319,-0.123) circle[radius=5.3pt];
\fill[g5,opacity=.6] (axis cs:0.327,-0.016) circle[radius=5.3pt];
\fill[g6,opacity=.6] (axis cs:0.996,0.466) circle[radius=5.3pt];
\fill[g1,opacity=.6] (axis cs:0.576,0.170) circle[radius=5.3pt];
\fill[g5,opacity=.6] (axis cs:0.853,-0.031) circle[radius=5.3pt];
\fill[g4,opacity=.6] (axis cs:-0.176,0.625) circle[radius=5.3pt];
\fill[g5,opacity=.6] (axis cs:0.659,0.192) circle[radius=5.0pt];
\fill[g6,opacity=.6] (axis cs:1.076,0.629) circle[radius=5.0pt];
\fill[g6,opacity=.6] (axis cs:1.252,1.231) circle[radius=5.0pt];
\fill[gray!50,opacity=.6] (axis cs:0.909,-1.615) circle[radius=5.0pt];
\fill[g3,opacity=.6] (axis cs:-0.067,-0.045) circle[radius=5.0pt];
\fill[g1,opacity=.6] (axis cs:0.694,-0.307) circle[radius=5.0pt];
\fill[g4,opacity=.6] (axis cs:-0.274,0.916) circle[radius=4.7pt];
\fill[g1,opacity=.6] (axis cs:0.005,0.259) circle[radius=4.7pt];
\fill[gray!50,opacity=.6] (axis cs:0.451,-1.043) circle[radius=4.7pt];
\fill[g1,opacity=.6] (axis cs:0.366,-0.319) circle[radius=4.7pt];
\fill[g7,opacity=.6] (axis cs:1.536,0.089) circle[radius=4.7pt];
\fill[g4,opacity=.6] (axis cs:-0.287,0.311) circle[radius=4.7pt];
\fill[g1,opacity=.6] (axis cs:0.427,-0.019) circle[radius=4.7pt];
\fill[g6,opacity=.6] (axis cs:1.611,1.435) circle[radius=4.4pt];
\fill[g1,opacity=.6] (axis cs:0.808,0.240) circle[radius=4.4pt];
\fill[g7,opacity=.6] (axis cs:1.004,-0.792) circle[radius=4.4pt];
\fill[gray!50,opacity=.6] (axis cs:0.469,-0.624) circle[radius=4.4pt];
\fill[g6,opacity=.6] (axis cs:0.535,0.584) circle[radius=4.4pt];
\fill[g2,opacity=.6] (axis cs:-0.169,0.610) circle[radius=4.4pt];
\fill[g2,opacity=.6] (axis cs:-0.033,-0.102) circle[radius=4.4pt];
\fill[g4,opacity=.6] (axis cs:0.061,0.409) circle[radius=4.4pt];
\fill[g4,opacity=.6] (axis cs:0.313,0.795) circle[radius=4.4pt];
\fill[g3,opacity=.6] (axis cs:0.346,-0.746) circle[radius=4.4pt];
\fill[g3,opacity=.6] (axis cs:0.917,0.130) circle[radius=4.0pt];
\fill[g2,opacity=.6] (axis cs:0.132,-0.104) circle[radius=4.0pt];
\fill[g2,opacity=.6] (axis cs:0.317,0.663) circle[radius=4.0pt];
\fill[g3,opacity=.6] (axis cs:0.184,0.341) circle[radius=4.0pt];
\fill[g5,opacity=.6] (axis cs:0.394,-0.347) circle[radius=4.0pt];
\fill[g3,opacity=.6] (axis cs:-0.347,-0.027) circle[radius=4.0pt];
\fill[g3,opacity=.6] (axis cs:-0.131,0.165) circle[radius=4.0pt];
\fill[g3,opacity=.6] (axis cs:-0.217,0.460) circle[radius=4.0pt];
\fill[g6,opacity=.6] (axis cs:1.341,0.984) circle[radius=4.0pt];
\fill[g6,opacity=.6] (axis cs:0.792,0.915) circle[radius=4.0pt];
\fill[g6,opacity=.6] (axis cs:1.257,0.483) circle[radius=4.0pt];
\fill[g8,opacity=.6] (axis cs:1.374,1.404) circle[radius=4.0pt];
\fill[g8,opacity=.6] (axis cs:0.882,0.671) circle[radius=4.0pt];
\fill[g6,opacity=.6] (axis cs:1.043,-0.212) circle[radius=4.0pt];
\fill[g3,opacity=.6] (axis cs:-0.157,-0.248) circle[radius=3.6pt];
\fill[g2,opacity=.6] (axis cs:-0.132,0.301) circle[radius=3.6pt];
\fill[g3,opacity=.6] (axis cs:0.176,0.449) circle[radius=3.6pt];
\fill[g6,opacity=.6] (axis cs:0.427,0.884) circle[radius=3.6pt];
\fill[g8,opacity=.6] (axis cs:0.645,0.887) circle[radius=3.6pt];
\fill[gray!50,opacity=.6] (axis cs:0.692,-1.196) circle[radius=3.6pt];
\fill[g6,opacity=.6] (axis cs:0.743,0.354) circle[radius=3.6pt];
\fill[gray!50,opacity=.6] (axis cs:1.019,-0.476) circle[radius=3.6pt];
\fill[gray!50,opacity=.6] (axis cs:1.477,-1.048) circle[radius=3.6pt];
\fill[g6,opacity=.6] (axis cs:0.776,0.077) circle[radius=3.6pt];
\fill[g8,opacity=.6] (axis cs:0.972,0.341) circle[radius=3.6pt];
\fill[g6,opacity=.6] (axis cs:1.348,-0.068) circle[radius=3.6pt];
\fill[g8,opacity=.6] (axis cs:1.849,1.655) circle[radius=3.6pt];
\fill[gray!50,opacity=.6] (axis cs:1.888,-0.360) circle[radius=3.6pt];
\fill[g6,opacity=.6] (axis cs:1.905,1.778) circle[radius=3.6pt];
\fill[g7,opacity=.6] (axis cs:2.011,1.244) circle[radius=3.6pt];
\fill[g8,opacity=.6] (axis cs:1.239,0.052) circle[radius=3.6pt];
\fill[g6,opacity=.6] (axis cs:1.424,0.935) circle[radius=3.6pt];
\fill[g2,opacity=.6] (axis cs:-0.217,-0.200) circle[radius=3.1pt];
\fill[gray!50,opacity=.6] (axis cs:1.331,-0.991) circle[radius=3.1pt];
\fill[gray!50,opacity=.6] (axis cs:1.258,-1.520) circle[radius=3.1pt];
\fill[g7,opacity=.6] (axis cs:1.790,0.500) circle[radius=3.1pt];
\fill[g6,opacity=.6] (axis cs:1.744,1.144) circle[radius=3.1pt];
\fill[g5,opacity=.6] (axis cs:-0.206,0.074) circle[radius=3.1pt];
\fill[g8,opacity=.6] (axis cs:1.194,1.075) circle[radius=3.1pt];
\fill[gray!50,opacity=.6] (axis cs:1.809,-0.645) circle[radius=3.1pt];
\fill[g7,opacity=.6] (axis cs:1.104,-0.463) circle[radius=3.1pt];
\fill[gray!50,opacity=.6] (axis cs:1.057,-0.791) circle[radius=3.1pt];
\fill[g1,opacity=.6] (axis cs:0.932,0.616) circle[radius=3.1pt];
\fill[g7,opacity=.6] (axis cs:0.882,-0.457) circle[radius=3.1pt];
\fill[g1,opacity=.6] (axis cs:0.797,0.167) circle[radius=3.1pt];
\fill[gray!50,opacity=.6] (axis cs:0.872,-1.307) circle[radius=3.1pt];
\fill[g2,opacity=.6] (axis cs:0.080,0.879) circle[radius=3.1pt];
\fill[g6,opacity=.6] (axis cs:0.068,0.394) circle[radius=3.1pt];
\fill[gray!50,opacity=.6] (axis cs:0.184,-0.754) circle[radius=3.1pt];
\fill[g5,opacity=.6] (axis cs:0.259,0.252) circle[radius=3.1pt];
\fill[g1,opacity=.6] (axis cs:0.192,0.433) circle[radius=3.1pt];
\fill[gray!50,opacity=.6] (axis cs:0.414,-0.735) circle[radius=3.1pt];
\fill[g3,opacity=.6] (axis cs:0.048,0.252) circle[radius=3.1pt];
\fill[g1,opacity=.6] (axis cs:0.443,0.820) circle[radius=3.1pt];
\fill[g3,opacity=.6] (axis cs:0.656,-0.330) circle[radius=3.1pt];
\fill[gray!50,opacity=.6] (axis cs:0.649,-0.451) circle[radius=3.1pt];
\fill[g8,opacity=.6] (axis cs:1.223,0.727) circle[radius=2.5pt];
\fill[g8,opacity=.6] (axis cs:0.845,0.980) circle[radius=2.5pt];
\fill[g6,opacity=.6] (axis cs:0.772,1.111) circle[radius=2.5pt];
\fill[g1,opacity=.6] (axis cs:0.970,0.307) circle[radius=2.5pt];
\fill[g3,opacity=.6] (axis cs:0.602,0.355) circle[radius=2.5pt];
\fill[g6,opacity=.6] (axis cs:1.357,0.309) circle[radius=2.5pt];
\fill[gray!50,opacity=.6] (axis cs:1.218,-0.391) circle[radius=2.5pt];
\fill[g8,opacity=.6] (axis cs:1.391,1.084) circle[radius=2.5pt];
\fill[gray!50,opacity=.6] (axis cs:1.970,0.650) circle[radius=2.5pt];
\fill[g1,opacity=.6] (axis cs:0.458,0.145) circle[radius=2.5pt];
\fill[g4,opacity=.6] (axis cs:-0.139,0.316) circle[radius=2.5pt];
\fill[gray!50,opacity=.6] (axis cs:0.147,-0.446) circle[radius=2.5pt];
\fill[g3,opacity=.6] (axis cs:-0.053,0.557) circle[radius=2.5pt];
\fill[g3,opacity=.6] (axis cs:0.333,0.166) circle[radius=2.5pt];
\fill[g4,opacity=.6] (axis cs:0.365,-0.188) circle[radius=2.5pt];
\fill[g5,opacity=.6] (axis cs:0.112,0.161) circle[radius=2.5pt];
\fill[g5,opacity=.6] (axis cs:-0.136,-0.291) circle[radius=2.5pt];
\fill[g5,opacity=.6] (axis cs:-0.210,0.206) circle[radius=2.5pt];
\fill[g8,opacity=.6] (axis cs:1.201,0.360) circle[radius=2.5pt];
\fill[g3,opacity=.6] (axis cs:0.093,-0.491) circle[radius=2.5pt];
\fill[g4,opacity=.6] (axis cs:0.099,0.101) circle[radius=2.5pt];
\fill[g6,opacity=.6] (axis cs:0.356,0.472) circle[radius=2.5pt];
\fill[g3,opacity=.6] (axis cs:0.524,-0.211) circle[radius=2.5pt];
\fill[g7,opacity=.6] (axis cs:0.669,-0.249) circle[radius=2.5pt];
\fill[g3,opacity=.6] (axis cs:-0.485,0.202) circle[radius=2.1pt];
\fill[gray!50,opacity=.6] (axis cs:0.496,-0.350) circle[radius=2.1pt];
\fill[g6,opacity=.6] (axis cs:0.315,0.912) circle[radius=2.1pt];
\fill[g6,opacity=.6] (axis cs:0.282,0.186) circle[radius=2.1pt];
\fill[gray!50,opacity=.6] (axis cs:0.241,-0.772) circle[radius=2.1pt];
\fill[g6,opacity=.6] (axis cs:0.105,0.085) circle[radius=2.1pt];
\fill[g3,opacity=.6] (axis cs:0.098,0.185) circle[radius=2.1pt];
\fill[g3,opacity=.6] (axis cs:0.307,-0.425) circle[radius=2.1pt];
\fill[gray!50,opacity=.6] (axis cs:0.186,0.057) circle[radius=2.1pt];
\fill[g5,opacity=.6] (axis cs:0.045,0.460) circle[radius=2.1pt];
\fill[g5,opacity=.6] (axis cs:0.016,0.606) circle[radius=2.1pt];
\fill[g5,opacity=.6] (axis cs:0.154,-0.553) circle[radius=2.1pt];
\fill[g4,opacity=.6] (axis cs:0.001,0.391) circle[radius=2.1pt];
\fill[g3,opacity=.6] (axis cs:-0.194,0.060) circle[radius=2.1pt];
\fill[g5,opacity=.6] (axis cs:-0.095,0.385) circle[radius=2.1pt];
\fill[g4,opacity=.6] (axis cs:-0.173,0.742) circle[radius=2.1pt];
\fill[gray!50,opacity=.6] (axis cs:-0.053,-0.539) circle[radius=2.1pt];
\fill[g5,opacity=.6] (axis cs:-0.410,0.113) circle[radius=2.1pt];
\fill[gray!50,opacity=.6] (axis cs:0.605,-1.018) circle[radius=2.1pt];
\fill[g6,opacity=.6] (axis cs:0.417,0.489) circle[radius=2.1pt];
\fill[g6,opacity=.6] (axis cs:0.586,-0.411) circle[radius=2.1pt];
\fill[g5,opacity=.6] (axis cs:0.431,0.247) circle[radius=2.1pt];
\fill[g3,opacity=.6] (axis cs:0.446,0.801) circle[radius=2.1pt];
\fill[g1,opacity=.6] (axis cs:0.845,-0.068) circle[radius=2.1pt];
\fill[g1,opacity=.6] (axis cs:0.756,0.515) circle[radius=2.1pt];
\fill[g7,opacity=.6] (axis cs:0.845,-0.149) circle[radius=2.1pt];
\fill[g3,opacity=.6] (axis cs:0.811,0.221) circle[radius=2.1pt];
\fill[g6,opacity=.6] (axis cs:0.856,0.086) circle[radius=2.1pt];
\fill[g3,opacity=.6] (axis cs:0.779,0.058) circle[radius=2.1pt];
\fill[g5,opacity=.6] (axis cs:0.645,0.039) circle[radius=2.1pt];
\fill[g4,opacity=.6] (axis cs:0.629,-0.273) circle[radius=2.1pt];
\fill[gray!50,opacity=.6] (axis cs:0.873,-0.418) circle[radius=2.1pt];
\fill[gray!50,opacity=.6] (axis cs:0.779,-0.239) circle[radius=2.1pt];
\fill[g6,opacity=.6] (axis cs:0.902,1.136) circle[radius=2.1pt];
\fill[g3,opacity=.6] (axis cs:0.902,0.805) circle[radius=2.1pt];
\fill[g8,opacity=.6] (axis cs:0.991,0.549) circle[radius=2.1pt];
\fill[g6,opacity=.6] (axis cs:0.908,0.387) circle[radius=2.1pt];
\fill[g6,opacity=.6] (axis cs:0.945,0.079) circle[radius=2.1pt];
\fill[gray!50,opacity=.6] (axis cs:0.924,0.056) circle[radius=2.1pt];
\fill[g6,opacity=.6] (axis cs:1.394,0.000) circle[radius=2.1pt];
\fill[g7,opacity=.6] (axis cs:1.411,-0.286) circle[radius=2.1pt];
\fill[gray!50,opacity=.6] (axis cs:1.207,0.352) circle[radius=2.1pt];
\fill[g8,opacity=.6] (axis cs:1.114,0.724) circle[radius=2.1pt];
\fill[gray!50,opacity=.6] (axis cs:1.221,-1.211) circle[radius=2.1pt];
\fill[g7,opacity=.6] (axis cs:1.166,-0.621) circle[radius=2.1pt];
\fill[g6,opacity=.6] (axis cs:1.178,1.140) circle[radius=2.1pt];
\fill[g6,opacity=.6] (axis cs:1.205,0.181) circle[radius=2.1pt];
\fill[g6,opacity=.6] (axis cs:1.329,1.124) circle[radius=2.1pt];
\fill[g8,opacity=.6] (axis cs:1.239,1.100) circle[radius=2.1pt];
\fill[g6,opacity=.6] (axis cs:1.429,0.623) circle[radius=2.1pt];
\fill[g8,opacity=.6] (axis cs:1.490,0.438) circle[radius=2.1pt];
\fill[g7,opacity=.6] (axis cs:1.498,0.398) circle[radius=2.1pt];
\fill[gray!50,opacity=.6] (axis cs:1.485,0.309) circle[radius=2.1pt];
\fill[g6,opacity=.6] (axis cs:1.389,0.748) circle[radius=2.1pt];
\fill[g8,opacity=.6] (axis cs:1.463,0.820) circle[radius=2.1pt];
\fill[g8,opacity=.6] (axis cs:1.370,1.536) circle[radius=2.1pt];
\fill[gray!50,opacity=.6] (axis cs:1.629,-0.534) circle[radius=2.1pt];
\fill[gray!50,opacity=.6] (axis cs:1.795,0.263) circle[radius=2.1pt];
\fill[g7,opacity=.6] (axis cs:1.642,-0.369) circle[radius=2.1pt];
\fill[g6,opacity=.6] (axis cs:1.680,1.336) circle[radius=2.1pt];
\fill[g6,opacity=.6] (axis cs:1.680,1.010) circle[radius=2.1pt];
\fill[g6,opacity=.6] (axis cs:1.995,1.530) circle[radius=2.1pt];
\fill[gray!50,opacity=.6] (axis cs:1.932,0.959) circle[radius=2.1pt];
\fill[g8,opacity=.6] (axis cs:2.101,2.042) circle[radius=2.1pt];
\fill[gray!50,opacity=.6] (axis cs:2.713,2.108) circle[radius=2.1pt];
\node[anchor=south west,font=\small\bfseries,text=tinta] at (axis cs:-0.85,1.15) {G2 · G4};
\node[anchor=south west,font=\small\bfseries,text=tinta] at (axis cs:-0.3,-1.55) {G1 · G3 · G5};
\node[anchor=south west,font=\small\bfseries,text=g6] at (axis cs:1.25,1.75) {G6};
\node[anchor=south west,font=\small\bfseries,text=g8] at (axis cs:2.35,1.3) {G8};
\node[anchor=south west,font=\small\bfseries,text=g7] at (axis cs:1.55,-1.05) {G7};
\end{axis}
\end{tikzpicture}

          \caption{Plano del ACM: ocho grupos estables (burbuja: igual perfil)}
        \end{figure}
      \end{column}
      \begin{column}{0.48\linewidth}
        \begin{figure}
          \centering
          % Generado por gen_tikz.py
\begin{tikzpicture}
\filldraw[fill=g1,draw=white,line width=0.8pt] (2.84,-6.29) -- (2.77,-6.26) -- (2.75,-6.17) -- (2.66,-6.10) -- (2.70,-6.05) -- (2.66,-5.96) -- (2.70,-5.87) -- (2.66,-5.82) -- (2.59,-5.82) -- (2.56,-5.66) -- (2.52,-5.64) -- (2.33,-5.68) -- (2.22,-5.66) -- (2.12,-5.68) -- (2.04,-5.53) -- (2.06,-5.39) -- (1.96,-5.32) -- (1.99,-5.27) -- (1.94,-5.27) -- (1.96,-5.24) -- (1.92,-5.25) -- (1.89,-5.22) -- (1.92,-5.20) -- (1.89,-5.09) -- (1.97,-5.06) -- (2.01,-5.09) -- (2.04,-5.04) -- (2.01,-4.94) -- (2.15,-4.90) -- (2.24,-4.95) -- (2.31,-4.81) -- (2.27,-4.71) -- (2.20,-4.65) -- (2.17,-4.53) -- (2.03,-4.42) -- (1.90,-4.41) -- (1.75,-4.20) -- (1.83,-4.16) -- (1.87,-4.02) -- (1.90,-3.98) -- (1.90,-3.90) -- (1.87,-3.82) -- (1.90,-3.82) -- (1.90,-3.88) -- (1.90,-3.86) -- (1.94,-3.88) -- (1.90,-3.91) -- (1.92,-3.97) -- (1.96,-3.93) -- (2.01,-3.91) -- (2.03,-3.93) -- (1.97,-3.98) -- (2.01,-4.02) -- (1.92,-4.02) -- (1.92,-4.12) -- (2.08,-4.11) -- (2.12,-3.98) -- (2.06,-3.93) -- (2.08,-3.68) -- (2.01,-3.60) -- (1.92,-3.54) -- (1.96,-3.49) -- (2.17,-3.45) -- (2.19,-3.40) -- (2.38,-3.26) -- (2.41,-3.33) -- (2.38,-3.37) -- (2.43,-3.47) -- (2.50,-3.47) -- (2.59,-3.68) -- (2.49,-3.77) -- (2.43,-3.88) -- (2.45,-3.91) -- (2.36,-3.98) -- (2.38,-4.12) -- (2.31,-4.27) -- (2.29,-4.41) -- (2.40,-4.58) -- (2.89,-4.62) -- (2.96,-4.41) -- (3.08,-4.42) -- (3.14,-4.37) -- (3.23,-4.20) -- (3.33,-4.16) -- (3.33,-4.09) -- (3.42,-4.04) -- (3.44,-3.98) -- (3.60,-4.00) -- (3.67,-3.97) -- (3.70,-3.98) -- (3.79,-3.86) -- (3.84,-3.93) -- (3.98,-4.02) -- (4.05,-4.12) -- (4.05,-4.23) -- (4.11,-4.34) -- (4.02,-4.37) -- (4.00,-4.53) -- (4.09,-4.64) -- (4.16,-4.53) -- (4.20,-4.51) -- (4.23,-4.55) -- (4.16,-4.64) -- (4.16,-4.76) -- (4.21,-4.92) -- (4.30,-4.94) -- (4.60,-4.67) -- (4.58,-4.78) -- (4.62,-4.94) -- (4.51,-5.01) -- (4.42,-5.15) -- (4.16,-5.31) -- (4.20,-5.46) -- (4.02,-5.64) -- (3.98,-5.71) -- (4.02,-5.76) -- (4.02,-5.83) -- (3.98,-5.85) -- (4.00,-5.90) -- (3.95,-5.94) -- (3.93,-6.03) -- (3.90,-6.03) -- (3.88,-6.08) -- (3.81,-6.10) -- (3.75,-6.05) -- (3.54,-6.13) -- (3.56,-6.19) -- (3.38,-6.35) -- (3.30,-6.12) -- (3.26,-6.10) -- (3.21,-6.12) -- (3.10,-6.06) -- (3.12,-6.26) -- (3.01,-6.26) -- (2.98,-6.22) -- cycle;
\filldraw[fill=g1,draw=white,line width=0.8pt] (3.88,-1.73) -- (3.72,-2.04) -- (3.56,-1.90) -- (3.42,-1.83) -- (3.42,-1.62) -- (3.45,-1.53) -- (3.60,-1.46) -- (3.61,-1.43) -- (3.60,-1.43) -- (3.70,-1.34) -- (3.77,-1.32) -- (3.77,-1.29) -- (3.88,-1.45) -- cycle;
\filldraw[fill=g7,draw=white,line width=0.8pt] (4.46,-6.87) -- (4.51,-6.87) -- (4.49,-7.03) -- (4.53,-7.05) -- (4.42,-7.23) -- (4.44,-7.30) -- (4.39,-7.39) -- (4.42,-7.44) -- (4.41,-7.56) -- (4.39,-7.60) -- (4.34,-7.60) -- (4.25,-7.72) -- (4.27,-7.77) -- (4.12,-7.84) -- (4.20,-7.69) -- (4.23,-7.49) -- (4.30,-7.51) -- (4.34,-7.39) -- (4.37,-7.12) -- (4.32,-7.05) -- (4.39,-7.02) -- (4.37,-6.98) -- cycle;
\filldraw[fill=g1,draw=white,line width=0.8pt] (4.04,-4.09) -- (3.84,-3.93) -- (3.79,-3.84) -- (3.84,-3.79) -- (3.95,-3.79) -- (3.98,-3.68) -- (4.02,-3.70) -- (4.04,-3.67) -- (3.98,-3.37) -- (4.04,-3.33) -- (3.95,-3.17) -- (3.79,-3.10) -- (3.72,-3.01) -- (3.68,-2.84) -- (3.74,-2.77) -- (3.61,-2.73) -- (3.63,-2.68) -- (3.47,-2.63) -- (3.49,-2.57) -- (3.38,-2.54) -- (3.35,-2.61) -- (3.30,-2.59) -- (3.33,-2.41) -- (3.24,-2.34) -- (3.23,-2.36) -- (3.21,-2.24) -- (3.16,-2.22) -- (3.19,-2.13) -- (3.12,-2.19) -- (3.17,-2.04) -- (3.01,-2.13) -- (3.10,-2.04) -- (3.12,-2.06) -- (3.08,-2.04) -- (3.12,-1.99) -- (3.14,-2.03) -- (3.19,-1.99) -- (3.17,-1.90) -- (3.14,-1.90) -- (3.19,-1.83) -- (3.17,-1.76) -- (3.28,-1.66) -- (3.38,-1.62) -- (3.40,-1.55) -- (3.44,-1.57) -- (3.42,-1.83) -- (3.67,-1.97) -- (3.72,-2.04) -- (3.68,-2.15) -- (3.82,-2.26) -- (3.75,-2.33) -- (3.74,-2.38) -- (3.82,-2.54) -- (3.84,-2.63) -- (3.81,-2.66) -- (3.84,-2.71) -- (3.81,-2.75) -- (3.84,-2.73) -- (3.93,-2.78) -- (4.07,-2.94) -- (4.14,-2.91) -- (4.16,-2.96) -- (4.21,-2.94) -- (4.27,-2.98) -- (4.27,-3.01) -- (4.30,-3.01) -- (4.42,-3.10) -- (4.41,-3.12) -- (4.49,-3.12) -- (4.55,-3.17) -- (4.62,-3.17) -- (4.71,-3.31) -- (4.69,-3.45) -- (4.74,-3.63) -- (4.74,-3.81) -- (4.72,-3.91) -- (4.65,-3.97) -- (4.64,-4.04) -- (4.71,-4.35) -- (4.65,-4.46) -- (4.60,-4.49) -- (4.60,-4.67) -- (4.30,-4.94) -- (4.21,-4.92) -- (4.16,-4.76) -- (4.18,-4.62) -- (4.23,-4.53) -- (4.20,-4.51) -- (4.16,-4.53) -- (4.09,-4.64) -- (3.98,-4.53) -- (4.02,-4.35) -- (4.11,-4.30) -- (4.05,-4.23) -- cycle;
\filldraw[fill=g1,draw=white,line width=0.8pt] (5.06,-6.82) -- (4.95,-6.82) -- (4.92,-6.79) -- (4.99,-6.59) -- (4.95,-6.57) -- (4.94,-6.47) -- (4.88,-6.43) -- (4.88,-6.36) -- (4.72,-6.27) -- (4.69,-6.20) -- (4.60,-6.24) -- (4.62,-6.31) -- (4.51,-6.40) -- (4.44,-6.29) -- (4.41,-6.33) -- (4.30,-6.27) -- (4.32,-6.24) -- (4.25,-6.17) -- (4.28,-6.10) -- (4.23,-5.98) -- (4.14,-6.05) -- (4.05,-6.01) -- (3.95,-6.06) -- (3.95,-5.96) -- (4.00,-5.90) -- (3.98,-5.85) -- (4.02,-5.82) -- (4.02,-5.76) -- (3.98,-5.71) -- (4.04,-5.59) -- (4.18,-5.78) -- (4.27,-5.75) -- (4.30,-5.76) -- (4.32,-5.85) -- (4.28,-5.98) -- (4.35,-5.96) -- (4.37,-5.90) -- (4.49,-5.99) -- (4.55,-6.08) -- (4.58,-6.06) -- (4.69,-6.08) -- (4.78,-6.03) -- (4.83,-6.08) -- (4.88,-5.94) -- (4.85,-5.90) -- (4.97,-5.73) -- (5.04,-5.76) -- (5.08,-5.85) -- (5.04,-5.92) -- (5.01,-5.92) -- (4.99,-5.99) -- (5.06,-6.05) -- (5.09,-5.94) -- (5.20,-5.99) -- (5.25,-5.96) -- (5.22,-5.92) -- (5.24,-5.82) -- (5.32,-5.85) -- (5.38,-5.80) -- (5.39,-5.85) -- (5.48,-5.73) -- (5.64,-5.62) -- (5.66,-5.41) -- (5.59,-5.32) -- (5.62,-5.31) -- (5.73,-5.43) -- (5.82,-5.39) -- (5.82,-5.31) -- (5.87,-5.24) -- (5.85,-5.01) -- (5.96,-5.04) -- (6.05,-4.92) -- (6.10,-4.95) -- (6.15,-4.88) -- (6.35,-4.92) -- (6.26,-5.15) -- (6.20,-5.16) -- (6.19,-5.36) -- (6.10,-5.45) -- (6.06,-5.43) -- (5.92,-5.66) -- (5.98,-5.76) -- (6.01,-5.75) -- (5.98,-5.80) -- (6.01,-5.90) -- (5.92,-5.90) -- (5.89,-5.96) -- (6.03,-6.01) -- (6.10,-6.12) -- (6.06,-6.13) -- (6.03,-6.26) -- (5.92,-6.22) -- (5.78,-6.40) -- (5.68,-6.47) -- (5.59,-6.38) -- (5.53,-6.42) -- (5.52,-6.40) -- (5.50,-6.49) -- (5.43,-6.57) -- (5.50,-6.66) -- (5.43,-6.73) -- (5.38,-6.73) -- (5.34,-6.89) -- (5.38,-6.94) -- (5.32,-7.02) -- (5.22,-7.00) -- (5.22,-6.96) -- (5.15,-6.96) -- cycle;
\filldraw[fill=g1,draw=white,line width=0.8pt] (3.38,-6.59) -- (3.31,-6.77) -- (3.35,-6.82) -- (3.30,-6.89) -- (3.21,-6.79) -- (3.08,-6.79) -- (3.07,-6.73) -- (3.01,-6.77) -- (2.98,-6.68) -- (2.93,-6.77) -- (2.87,-6.77) -- (2.82,-6.68) -- (2.84,-6.59) -- (2.89,-6.61) -- (2.93,-6.54) -- (2.91,-6.47) -- (3.01,-6.49) -- (3.07,-6.43) -- (3.00,-6.36) -- (2.87,-6.38) -- (2.87,-6.27) -- (2.98,-6.22) -- (3.01,-6.26) -- (3.12,-6.26) -- (3.10,-6.06) -- (3.21,-6.12) -- (3.26,-6.10) -- (3.30,-6.12) -- (3.35,-6.31) -- (3.38,-6.29) -- (3.38,-6.35) -- (3.44,-6.27) -- (3.56,-6.19) -- (3.54,-6.13) -- (3.75,-6.05) -- (3.84,-6.10) -- (3.90,-6.03) -- (3.93,-6.03) -- (3.97,-6.10) -- (3.93,-6.26) -- (3.95,-6.33) -- (3.88,-6.47) -- (3.79,-6.43) -- (3.61,-6.45) -- (3.54,-6.57) -- cycle;
\filldraw[fill=g1,draw=white,line width=0.8pt] (3.26,-10.51) -- (3.17,-10.47) -- (3.07,-10.38) -- (2.96,-10.42) -- (2.93,-10.36) -- (2.87,-10.38) -- (2.78,-10.29) -- (2.82,-10.24) -- (2.52,-10.17) -- (2.47,-10.08) -- (2.56,-9.99) -- (2.61,-9.80) -- (2.59,-9.75) -- (2.64,-9.78) -- (2.71,-9.76) -- (2.82,-9.64) -- (2.98,-9.59) -- (3.16,-9.29) -- (3.24,-9.25) -- (3.30,-9.11) -- (3.37,-9.06) -- (3.44,-8.94) -- (3.47,-8.94) -- (3.53,-8.87) -- (3.51,-8.74) -- (3.58,-8.73) -- (3.68,-8.58) -- (3.72,-8.60) -- (3.81,-8.55) -- (3.84,-8.62) -- (3.93,-8.60) -- (3.93,-8.67) -- (3.97,-8.73) -- (3.95,-8.76) -- (4.00,-8.80) -- (3.91,-8.88) -- (3.95,-8.95) -- (3.86,-9.15) -- (3.91,-9.13) -- (3.93,-9.20) -- (4.00,-9.20) -- (4.00,-9.27) -- (4.09,-9.25) -- (4.11,-9.31) -- (4.30,-9.34) -- (4.35,-9.39) -- (4.49,-9.41) -- (4.83,-9.73) -- (4.92,-9.92) -- (5.02,-10.01) -- (5.06,-10.14) -- (5.20,-10.28) -- (5.25,-10.28) -- (5.29,-10.35) -- (5.41,-10.36) -- (5.46,-10.26) -- (5.53,-10.26) -- (5.53,-10.14) -- (5.62,-10.12) -- (5.64,-10.17) -- (5.66,-10.12) -- (5.71,-10.15) -- (5.73,-10.12) -- (5.85,-10.22) -- (5.89,-10.21) -- (5.96,-10.35) -- (5.99,-10.35) -- (6.01,-10.42) -- (6.05,-10.38) -- (6.03,-10.45) -- (6.06,-10.49) -- (6.08,-10.47) -- (6.10,-10.54) -- (6.19,-10.51) -- (6.24,-10.56) -- (6.22,-10.58) -- (6.31,-10.58) -- (6.29,-10.65) -- (6.35,-10.68) -- (6.36,-10.66) -- (6.36,-10.77) -- (6.43,-10.77) -- (6.45,-10.80) -- (6.43,-10.86) -- (6.50,-10.84) -- (6.49,-10.88) -- (6.54,-10.91) -- (6.56,-10.89) -- (6.54,-10.91) -- (6.63,-11.00) -- (6.70,-10.98) -- (6.66,-11.02) -- (6.72,-11.02) -- (6.72,-11.07) -- (6.77,-11.00) -- (6.79,-11.05) -- (6.84,-11.00) -- (6.96,-11.12) -- (6.79,-11.18) -- (6.68,-11.16) -- (6.56,-11.23) -- (6.54,-11.32) -- (6.56,-11.37) -- (6.26,-11.40) -- (6.24,-11.46) -- (6.17,-11.49) -- (6.15,-11.56) -- (6.10,-11.62) -- (6.08,-11.70) -- (6.03,-11.76) -- (5.92,-11.67) -- (5.80,-11.79) -- (5.64,-11.67) -- (5.48,-11.72) -- (5.39,-11.63) -- (5.34,-11.72) -- (5.24,-11.72) -- (5.11,-11.62) -- (5.08,-11.65) -- (4.95,-11.63) -- (4.94,-11.60) -- (4.90,-11.63) -- (4.90,-11.56) -- (4.83,-11.58) -- (4.85,-11.55) -- (4.78,-11.51) -- (4.69,-11.53) -- (4.51,-11.47) -- (4.42,-11.39) -- (4.35,-11.40) -- (4.35,-11.35) -- (4.32,-11.39) -- (4.27,-11.35) -- (4.28,-11.30) -- (4.21,-11.26) -- (4.12,-11.30) -- (4.14,-11.26) -- (4.12,-11.23) -- (4.09,-11.30) -- (4.05,-11.25) -- (4.05,-11.28) -- (3.97,-11.25) -- (3.97,-11.21) -- (3.91,-11.23) -- (3.95,-11.18) -- (3.90,-11.14) -- (3.95,-11.10) -- (3.90,-11.09) -- (3.93,-11.09) -- (3.91,-11.02) -- (3.63,-10.93) -- (3.65,-10.89) -- (3.61,-10.82) -- (3.65,-10.80) -- (3.61,-10.75) -- (3.54,-10.75) -- (3.54,-10.72) -- (3.44,-10.73) -- (3.42,-10.66) -- (3.45,-10.63) -- (3.38,-10.58) -- (3.40,-10.52) -- (3.33,-10.49) -- cycle;
\filldraw[fill=g2,draw=white,line width=0.8pt] (1.92,-9.62) -- (1.73,-9.69) -- (1.73,-9.68) -- (1.64,-9.69) -- (1.57,-9.68) -- (1.60,-9.52) -- (1.67,-9.43) -- (1.57,-9.34) -- (1.59,-9.22) -- (1.53,-9.15) -- (1.45,-9.13) -- (1.43,-9.22) -- (1.20,-9.27) -- (1.11,-9.24) -- (1.04,-9.02) -- (1.06,-8.78) -- (1.16,-8.78) -- (1.15,-8.71) -- (1.20,-8.71) -- (1.18,-8.67) -- (1.27,-8.55) -- (1.27,-8.51) -- (1.23,-8.51) -- (1.37,-8.30) -- (1.41,-8.34) -- (1.48,-8.27) -- (1.53,-8.37) -- (1.57,-8.34) -- (1.66,-8.41) -- (1.76,-8.34) -- (1.92,-8.39) -- (1.94,-8.32) -- (2.06,-8.28) -- (2.10,-8.39) -- (2.26,-8.41) -- (2.29,-8.39) -- (2.34,-8.20) -- (2.54,-8.25) -- (2.57,-8.30) -- (2.70,-8.30) -- (2.64,-8.41) -- (2.73,-8.50) -- (2.71,-8.55) -- (2.86,-8.65) -- (2.91,-8.64) -- (2.93,-8.71) -- (2.96,-8.73) -- (2.89,-8.88) -- (2.91,-8.95) -- (2.84,-9.01) -- (2.77,-8.94) -- (2.71,-9.02) -- (2.66,-9.02) -- (2.61,-9.13) -- (2.64,-9.18) -- (2.61,-9.22) -- (2.56,-9.22) -- (2.49,-9.13) -- (2.41,-9.15) -- (2.36,-9.22) -- (2.36,-9.29) -- (2.33,-9.31) -- (2.26,-9.27) -- (2.22,-9.29) -- (2.24,-9.38) -- (2.19,-9.41) -- (2.22,-9.50) -- (2.36,-9.57) -- (2.38,-9.64) -- (2.41,-9.62) -- (2.41,-9.68) -- (2.54,-9.68) -- (2.57,-9.76) -- (2.61,-9.76) -- (2.56,-9.99) -- (2.49,-10.06) -- (2.49,-10.12) -- (2.52,-10.17) -- (2.64,-10.19) -- (2.68,-10.26) -- (2.59,-10.26) -- (2.54,-10.31) -- (2.36,-10.31) -- (2.26,-10.24) -- (2.24,-9.94) -- (2.15,-9.87) -- (1.94,-9.98) -- (1.92,-9.84) -- (1.99,-9.82) -- (2.01,-9.73) -- cycle;
\filldraw[fill=g1,draw=white,line width=0.8pt] (5.13,-4.04) -- (5.16,-4.11) -- (5.09,-4.14) -- (5.11,-4.21) -- (5.02,-4.32) -- (4.97,-4.34) -- (4.88,-4.28) -- (4.74,-4.27) -- (4.74,-4.21) -- (4.81,-4.12) -- (4.71,-3.95) -- (4.74,-3.63) -- (4.69,-3.47) -- (4.71,-3.31) -- (4.64,-3.23) -- (4.71,-3.10) -- (4.64,-3.03) -- (4.64,-3.00) -- (4.57,-2.98) -- (4.58,-2.89) -- (4.53,-2.87) -- (4.55,-2.80) -- (4.41,-2.71) -- (4.49,-2.63) -- (4.71,-2.59) -- (4.67,-2.45) -- (4.44,-2.20) -- (4.46,-2.12) -- (4.60,-1.92) -- (4.72,-1.90) -- (4.85,-1.82) -- (4.92,-1.82) -- (4.86,-1.76) -- (4.86,-1.69) -- (4.94,-1.59) -- (4.85,-1.55) -- (4.99,-1.50) -- (5.15,-1.50) -- (5.18,-1.53) -- (5.16,-1.62) -- (5.20,-1.60) -- (5.32,-1.66) -- (5.32,-1.76) -- (5.25,-1.85) -- (5.29,-1.90) -- (5.50,-1.87) -- (5.52,-1.92) -- (5.48,-2.17) -- (5.41,-2.33) -- (5.46,-2.40) -- (5.36,-2.56) -- (5.34,-2.66) -- (5.29,-2.66) -- (5.11,-2.96) -- (5.04,-3.01) -- (5.01,-3.28) -- (5.04,-3.37) -- (4.92,-3.51) -- (4.95,-3.58) -- (4.90,-3.60) -- (4.99,-3.63) -- (4.94,-3.70) -- (5.02,-3.74) -- (5.06,-3.70) -- (5.04,-3.63) -- (5.11,-3.65) -- (5.08,-3.81) -- (5.02,-3.81) -- (5.01,-3.88) -- (5.02,-3.93) -- (5.08,-3.95) -- (5.09,-4.04) -- cycle;
\filldraw[fill=g2,draw=white,line width=0.8pt] (3.10,-4.41) -- (2.96,-4.41) -- (2.89,-4.62) -- (2.40,-4.58) -- (2.29,-4.41) -- (2.33,-4.20) -- (2.38,-4.12) -- (2.36,-3.98) -- (2.45,-3.91) -- (2.43,-3.88) -- (2.49,-3.77) -- (2.59,-3.68) -- (2.50,-3.47) -- (2.43,-3.47) -- (2.38,-3.37) -- (2.41,-3.33) -- (2.38,-3.26) -- (2.50,-3.16) -- (2.57,-3.05) -- (2.59,-2.94) -- (2.80,-2.75) -- (2.93,-2.77) -- (2.89,-2.80) -- (3.01,-2.78) -- (3.03,-2.86) -- (3.10,-2.86) -- (3.10,-2.91) -- (3.17,-2.93) -- (3.19,-2.91) -- (3.24,-3.03) -- (3.30,-3.01) -- (3.38,-3.08) -- (3.44,-3.08) -- (3.49,-3.21) -- (3.42,-3.26) -- (3.38,-3.23) -- (3.30,-3.31) -- (3.35,-3.37) -- (3.31,-3.44) -- (3.35,-3.60) -- (3.44,-3.61) -- (3.47,-3.68) -- (3.53,-3.63) -- (3.60,-3.63) -- (3.63,-3.58) -- (3.70,-3.60) -- (3.70,-3.65) -- (3.81,-3.70) -- (3.79,-3.72) -- (3.82,-3.79) -- (3.70,-3.98) -- (3.53,-4.00) -- (3.44,-3.98) -- (3.42,-4.04) -- (3.33,-4.09) -- (3.33,-4.16) -- (3.23,-4.20) -- cycle;
\filldraw[fill=g1,draw=white,line width=0.8pt] (4.11,-7.49) -- (4.02,-7.37) -- (3.95,-7.44) -- (3.82,-7.35) -- (3.74,-7.35) -- (3.82,-7.16) -- (3.79,-6.96) -- (3.86,-6.91) -- (3.86,-6.75) -- (3.90,-6.73) -- (3.86,-6.72) -- (3.88,-6.47) -- (3.95,-6.33) -- (3.93,-6.26) -- (3.97,-6.10) -- (3.93,-6.05) -- (4.00,-6.05) -- (4.07,-6.01) -- (4.14,-6.05) -- (4.23,-5.98) -- (4.28,-6.10) -- (4.25,-6.17) -- (4.30,-6.24) -- (4.30,-6.29) -- (4.34,-6.27) -- (4.41,-6.33) -- (4.44,-6.29) -- (4.46,-6.35) -- (4.53,-6.38) -- (4.62,-6.31) -- (4.60,-6.24) -- (4.69,-6.20) -- (4.72,-6.27) -- (4.88,-6.36) -- (4.88,-6.43) -- (4.94,-6.47) -- (4.95,-6.57) -- (4.99,-6.59) -- (4.94,-6.80) -- (5.06,-6.82) -- (5.15,-6.96) -- (5.22,-6.96) -- (5.22,-7.00) -- (5.32,-7.02) -- (5.36,-6.94) -- (5.29,-7.42) -- (5.09,-7.31) -- (4.99,-7.37) -- (4.92,-7.33) -- (4.88,-7.21) -- (4.90,-7.14) -- (4.86,-7.10) -- (4.69,-7.23) -- (4.76,-7.44) -- (4.58,-7.47) -- (4.57,-7.53) -- (4.51,-7.51) -- (4.41,-7.61) -- (4.42,-7.46) -- (4.39,-7.39) -- (4.44,-7.30) -- (4.42,-7.23) -- (4.53,-7.05) -- (4.49,-7.03) -- (4.51,-6.87) -- (4.46,-6.86) -- (4.41,-6.96) -- (4.37,-6.98) -- (4.39,-7.02) -- (4.32,-7.05) -- (4.37,-7.12) -- (4.34,-7.39) -- (4.30,-7.51) -- (4.23,-7.49) -- (4.20,-7.69) -- (4.12,-7.84) -- (4.12,-7.77) -- (4.04,-7.74) -- (4.05,-7.56) -- cycle;
\filldraw[fill=g2,draw=white,line width=0.8pt] (2.26,-7.24) -- (2.34,-7.46) -- (2.22,-7.63) -- (2.19,-7.60) -- (2.13,-7.63) -- (2.03,-7.60) -- (1.90,-7.49) -- (1.76,-7.51) -- (1.71,-7.44) -- (1.66,-7.47) -- (1.64,-7.40) -- (1.59,-7.46) -- (1.55,-7.44) -- (1.46,-7.47) -- (1.46,-7.58) -- (1.39,-7.49) -- (1.39,-7.37) -- (1.50,-7.35) -- (1.45,-7.33) -- (1.55,-7.23) -- (1.57,-6.94) -- (1.59,-6.94) -- (1.59,-6.91) -- (1.55,-6.93) -- (1.52,-6.77) -- (1.52,-6.33) -- (1.48,-6.33) -- (1.46,-6.27) -- (1.37,-6.27) -- (1.43,-6.17) -- (1.55,-6.17) -- (1.64,-6.03) -- (1.55,-5.78) -- (1.53,-5.78) -- (1.55,-5.83) -- (1.53,-5.83) -- (1.43,-5.66) -- (1.45,-5.55) -- (1.50,-5.62) -- (1.53,-5.52) -- (1.50,-5.48) -- (1.53,-5.46) -- (1.52,-5.45) -- (1.57,-5.34) -- (1.50,-5.27) -- (1.50,-5.22) -- (1.45,-5.20) -- (1.39,-5.25) -- (1.34,-5.09) -- (1.27,-5.04) -- (1.23,-5.08) -- (1.27,-4.94) -- (1.08,-4.74) -- (1.15,-4.53) -- (1.23,-4.48) -- (1.20,-4.39) -- (1.22,-4.32) -- (1.29,-4.35) -- (1.36,-4.48) -- (1.43,-4.42) -- (1.52,-4.27) -- (1.55,-4.25) -- (1.53,-4.20) -- (1.57,-4.14) -- (1.71,-4.11) -- (1.66,-4.02) -- (1.62,-3.86) -- (1.55,-3.77) -- (1.50,-3.61) -- (1.46,-3.63) -- (1.46,-3.54) -- (1.52,-3.53) -- (1.50,-3.47) -- (1.53,-3.47) -- (1.53,-3.44) -- (1.64,-3.63) -- (1.75,-3.68) -- (1.82,-3.81) -- (1.87,-3.82) -- (1.90,-3.90) -- (1.90,-3.98) -- (1.87,-4.02) -- (1.83,-4.16) -- (1.75,-4.20) -- (1.75,-4.23) -- (1.90,-4.41) -- (2.03,-4.42) -- (2.17,-4.53) -- (2.20,-4.65) -- (2.27,-4.71) -- (2.31,-4.81) -- (2.24,-4.95) -- (2.15,-4.90) -- (2.01,-4.94) -- (2.04,-5.06) -- (1.89,-5.11) -- (1.92,-5.20) -- (1.89,-5.22) -- (1.92,-5.25) -- (1.96,-5.24) -- (1.94,-5.27) -- (1.99,-5.27) -- (1.96,-5.32) -- (2.06,-5.39) -- (2.04,-5.57) -- (2.13,-5.68) -- (2.54,-5.64) -- (2.59,-5.82) -- (2.66,-5.82) -- (2.70,-5.87) -- (2.66,-5.96) -- (2.70,-6.05) -- (2.66,-6.10) -- (2.75,-6.17) -- (2.73,-6.20) -- (2.68,-6.22) -- (2.68,-6.31) -- (2.57,-6.35) -- (2.56,-6.40) -- (2.63,-6.56) -- (2.70,-6.59) -- (2.66,-6.77) -- (2.63,-6.86) -- (2.47,-6.94) -- (2.49,-6.98) -- (2.47,-7.07) -- (2.38,-7.14) -- (2.41,-7.17) -- cycle;
\filldraw[fill=g1,draw=white,line width=0.8pt] (3.16,-9.29) -- (2.98,-9.59) -- (2.84,-9.64) -- (2.77,-9.73) -- (2.64,-9.78) -- (2.57,-9.76) -- (2.54,-9.68) -- (2.41,-9.68) -- (2.41,-9.62) -- (2.38,-9.64) -- (2.36,-9.57) -- (2.22,-9.50) -- (2.19,-9.41) -- (2.24,-9.38) -- (2.22,-9.29) -- (2.26,-9.27) -- (2.33,-9.31) -- (2.36,-9.29) -- (2.36,-9.22) -- (2.41,-9.15) -- (2.49,-9.13) -- (2.56,-9.22) -- (2.61,-9.22) -- (2.64,-9.18) -- (2.61,-9.13) -- (2.66,-9.02) -- (2.71,-9.02) -- (2.77,-8.94) -- (2.84,-9.01) -- (2.91,-8.95) -- (2.89,-8.88) -- (2.96,-8.73) -- (2.93,-8.71) -- (2.91,-8.64) -- (2.86,-8.65) -- (2.71,-8.55) -- (2.75,-8.51) -- (2.78,-8.57) -- (2.89,-8.62) -- (3.01,-8.50) -- (3.07,-8.50) -- (3.17,-8.21) -- (3.31,-8.13) -- (3.35,-8.13) -- (3.40,-8.18) -- (3.44,-8.13) -- (3.47,-8.16) -- (3.51,-8.11) -- (3.60,-8.11) -- (3.58,-8.23) -- (3.72,-8.23) -- (3.84,-8.11) -- (3.82,-8.04) -- (3.88,-8.00) -- (3.86,-7.97) -- (4.04,-7.81) -- (4.05,-7.76) -- (4.12,-7.77) -- (4.11,-7.84) -- (4.16,-7.86) -- (3.97,-8.07) -- (3.98,-8.25) -- (3.93,-8.34) -- (3.58,-8.73) -- (3.51,-8.74) -- (3.53,-8.87) -- (3.47,-8.94) -- (3.44,-8.94) -- (3.37,-9.06) -- (3.30,-9.11) -- (3.24,-9.25) -- cycle;
\filldraw[fill=g1,draw=white,line width=0.8pt] (5.52,-1.87) -- (5.43,-1.90) -- (5.34,-1.89) -- (5.29,-1.90) -- (5.25,-1.85) -- (5.32,-1.78) -- (5.34,-1.71) -- (5.27,-1.62) -- (5.16,-1.62) -- (5.18,-1.52) -- (4.99,-1.50) -- (4.88,-1.52) -- (4.83,-1.37) -- (4.85,-1.34) -- (4.83,-1.25) -- (4.90,-1.20) -- (4.90,-1.15) -- (5.16,-1.11) -- (5.62,-0.78) -- (6.08,-0.58) -- (6.17,-0.39) -- (6.13,-0.28) -- (6.33,-0.25) -- (6.29,-0.32) -- (6.31,-0.35) -- (6.33,-0.32) -- (6.36,-0.35) -- (6.42,-0.28) -- (6.40,-0.23) -- (6.35,-0.25) -- (6.43,-0.16) -- (6.47,-0.16) -- (6.45,-0.18) -- (6.47,-0.21) -- (6.52,-0.18) -- (6.50,-0.14) -- (6.54,-0.12) -- (6.57,-0.16) -- (6.63,-0.12) -- (6.57,-0.09) -- (6.54,-0.12) -- (6.57,-0.07) -- (6.72,-0.09) -- (6.98,-0.21) -- (7.09,-0.41) -- (7.07,-0.48) -- (6.94,-0.56) -- (6.89,-0.55) -- (6.93,-0.56) -- (6.91,-0.63) -- (6.33,-0.79) -- (6.08,-1.25) -- (5.99,-1.23) -- (5.85,-1.30) -- (5.82,-1.43) -- (5.73,-1.48) -- cycle;
\filldraw[fill=g2,draw=white,line width=0.8pt] (4.16,-2.94) -- (4.14,-2.91) -- (4.05,-2.94) -- (3.97,-2.82) -- (3.81,-2.75) -- (3.84,-2.71) -- (3.81,-2.66) -- (3.84,-2.63) -- (3.82,-2.54) -- (3.74,-2.38) -- (3.75,-2.33) -- (3.82,-2.26) -- (3.68,-2.13) -- (3.75,-1.92) -- (3.88,-1.78) -- (3.88,-1.45) -- (3.77,-1.29) -- (4.04,-1.37) -- (4.21,-1.39) -- (4.34,-1.30) -- (4.30,-1.18) -- (4.39,-1.08) -- (4.55,-1.08) -- (4.71,-1.15) -- (4.90,-1.15) -- (4.88,-1.22) -- (4.83,-1.25) -- (4.85,-1.34) -- (4.83,-1.37) -- (4.88,-1.52) -- (4.85,-1.57) -- (4.94,-1.59) -- (4.86,-1.69) -- (4.86,-1.76) -- (4.92,-1.82) -- (4.85,-1.82) -- (4.72,-1.90) -- (4.60,-1.92) -- (4.49,-2.06) -- (4.44,-2.20) -- (4.67,-2.45) -- (4.71,-2.59) -- (4.49,-2.63) -- (4.41,-2.71) -- (4.55,-2.80) -- (4.53,-2.87) -- (4.58,-2.89) -- (4.57,-2.98) -- (4.64,-3.00) -- (4.64,-3.03) -- (4.71,-3.12) -- (4.64,-3.21) -- (4.49,-3.12) -- (4.41,-3.12) -- (4.42,-3.10) -- (4.30,-3.01) -- (4.27,-3.01) -- (4.27,-2.98) -- cycle;
\filldraw[fill=g1,draw=white,line width=0.8pt] (5.01,-9.04) -- (5.01,-9.08) -- (4.95,-9.10) -- (4.95,-9.04) -- (4.90,-9.08) -- (4.88,-9.04) -- (4.83,-9.13) -- (4.83,-9.73) -- (4.49,-9.41) -- (4.35,-9.39) -- (4.30,-9.34) -- (4.11,-9.31) -- (4.09,-9.25) -- (4.00,-9.27) -- (4.00,-9.20) -- (3.93,-9.20) -- (3.91,-9.13) -- (3.86,-9.15) -- (3.95,-8.95) -- (3.91,-8.88) -- (4.00,-8.80) -- (3.95,-8.76) -- (3.97,-8.73) -- (3.93,-8.67) -- (3.93,-8.60) -- (3.84,-8.62) -- (3.81,-8.55) -- (3.72,-8.60) -- (3.70,-8.58) -- (3.93,-8.34) -- (3.98,-8.25) -- (3.97,-8.07) -- (4.18,-7.81) -- (4.27,-7.77) -- (4.25,-7.72) -- (4.34,-7.60) -- (4.41,-7.61) -- (4.51,-7.51) -- (4.57,-7.53) -- (4.58,-7.47) -- (4.76,-7.44) -- (4.69,-7.23) -- (4.86,-7.10) -- (4.90,-7.14) -- (4.88,-7.21) -- (4.92,-7.33) -- (4.99,-7.37) -- (5.09,-7.31) -- (5.29,-7.42) -- (5.36,-6.94) -- (5.46,-7.14) -- (5.69,-7.35) -- (5.69,-7.33) -- (5.80,-7.33) -- (5.82,-7.30) -- (5.89,-7.33) -- (5.92,-7.30) -- (5.94,-7.33) -- (6.08,-7.23) -- (6.17,-7.21) -- (6.24,-7.28) -- (6.31,-7.26) -- (6.47,-7.09) -- (6.54,-7.10) -- (7.12,-6.80) -- (7.12,-8.60) -- (7.07,-8.58) -- (7.05,-8.62) -- (7.03,-8.58) -- (6.94,-8.62) -- (6.89,-8.57) -- (6.84,-8.64) -- (6.77,-8.58) -- (6.77,-8.64) -- (6.66,-8.62) -- (6.68,-8.64) -- (6.61,-8.65) -- (6.59,-8.62) -- (6.57,-8.64) -- (6.50,-8.58) -- (6.50,-8.65) -- (6.45,-8.60) -- (6.43,-8.62) -- (6.45,-8.64) -- (6.31,-8.67) -- (6.22,-8.65) -- (6.19,-8.60) -- (6.12,-8.62) -- (6.13,-8.71) -- (6.06,-8.65) -- (6.05,-8.69) -- (6.01,-8.67) -- (5.94,-8.74) -- (5.89,-8.74) -- (5.89,-8.78) -- (5.82,-8.76) -- (5.78,-8.87) -- (5.75,-8.83) -- (5.73,-8.88) -- (5.69,-8.83) -- (5.64,-8.88) -- (5.57,-8.85) -- (5.48,-8.88) -- (5.45,-8.94) -- (5.46,-8.97) -- (5.45,-8.95) -- (5.27,-9.06) -- (5.27,-9.02) -- (5.22,-9.02) -- (5.18,-9.04) -- (5.20,-9.08) -- (5.11,-9.06) -- (5.09,-9.10) -- (5.09,-9.06) -- (5.08,-9.10) -- (5.04,-9.08) -- (5.02,-9.10) -- cycle;
\filldraw[fill=g2,draw=white,line width=0.8pt] (1.25,-10.43) -- (1.04,-10.45) -- (0.81,-10.33) -- (0.72,-10.22) -- (0.74,-10.19) -- (0.69,-10.12) -- (0.58,-10.15) -- (0.58,-10.10) -- (0.51,-10.12) -- (0.49,-10.06) -- (0.42,-10.06) -- (0.23,-9.87) -- (0.26,-9.85) -- (0.21,-9.85) -- (0.19,-9.78) -- (0.18,-9.80) -- (0.07,-9.71) -- (0.25,-9.54) -- (0.28,-9.55) -- (0.25,-9.57) -- (0.33,-9.55) -- (0.26,-9.57) -- (0.32,-9.57) -- (0.32,-9.61) -- (0.28,-9.62) -- (0.33,-9.61) -- (0.33,-9.57) -- (0.35,-9.61) -- (0.39,-9.57) -- (0.44,-9.57) -- (0.46,-9.62) -- (0.48,-9.61) -- (0.49,-9.52) -- (0.48,-9.48) -- (0.44,-9.48) -- (0.46,-9.41) -- (0.41,-9.38) -- (0.37,-9.41) -- (0.39,-9.34) -- (0.35,-9.29) -- (0.35,-9.18) -- (0.48,-8.99) -- (0.48,-8.95) -- (0.56,-8.85) -- (0.67,-8.80) -- (0.83,-8.78) -- (0.92,-8.85) -- (0.99,-8.80) -- (1.04,-8.81) -- (1.04,-9.02) -- (1.11,-9.22) -- (1.16,-9.25) -- (1.43,-9.22) -- (1.45,-9.13) -- (1.53,-9.15) -- (1.59,-9.22) -- (1.57,-9.34) -- (1.67,-9.43) -- (1.60,-9.52) -- (1.57,-9.68) -- (1.64,-9.69) -- (1.73,-9.68) -- (1.73,-9.69) -- (1.92,-9.62) -- (2.01,-9.73) -- (1.99,-9.82) -- (1.92,-9.84) -- (1.94,-9.99) -- (1.76,-10.14) -- (1.78,-10.21) -- (1.83,-10.24) -- (1.83,-10.28) -- (1.75,-10.38) -- (1.76,-10.43) -- (1.73,-10.43) -- (1.69,-10.61) -- (1.75,-10.65) -- (1.76,-10.84) -- (1.39,-10.80) -- (1.43,-10.58) -- (1.36,-10.59) -- (1.27,-10.54) -- (1.23,-10.51) -- (1.27,-10.47) -- cycle;
\filldraw[fill=g1,draw=white,line width=0.8pt] (4.97,-4.34) -- (5.11,-4.21) -- (5.09,-4.14) -- (5.16,-4.09) -- (5.13,-4.04) -- (5.09,-4.04) -- (5.08,-3.95) -- (5.02,-3.93) -- (5.01,-3.88) -- (5.02,-3.81) -- (5.08,-3.81) -- (5.11,-3.67) -- (5.04,-3.63) -- (5.06,-3.70) -- (5.02,-3.74) -- (4.94,-3.70) -- (4.99,-3.63) -- (4.92,-3.63) -- (4.90,-3.60) -- (4.95,-3.58) -- (4.92,-3.51) -- (4.95,-3.44) -- (5.04,-3.37) -- (5.01,-3.28) -- (5.04,-3.05) -- (5.08,-3.00) -- (5.25,-3.00) -- (5.41,-2.89) -- (5.46,-3.07) -- (5.57,-3.03) -- (5.62,-3.07) -- (5.71,-3.49) -- (5.96,-3.74) -- (5.99,-4.04) -- (5.92,-4.02) -- (5.87,-4.09) -- (5.90,-4.18) -- (5.87,-4.49) -- (5.92,-4.58) -- (6.13,-4.60) -- (6.17,-4.64) -- (6.15,-4.71) -- (6.26,-4.90) -- (6.13,-4.88) -- (6.08,-4.95) -- (6.05,-4.92) -- (5.96,-5.04) -- (5.87,-5.01) -- (5.82,-5.04) -- (5.80,-4.94) -- (5.71,-4.97) -- (5.69,-4.94) -- (5.64,-4.95) -- (5.64,-4.88) -- (5.55,-4.85) -- (5.55,-4.76) -- (5.52,-4.71) -- (5.55,-4.62) -- (5.50,-4.51) -- (5.43,-4.46) -- (5.43,-4.39) -- (5.22,-4.37) -- (5.18,-4.46) -- (4.97,-4.41) -- (4.85,-4.27) -- cycle;
\filldraw[fill=g1,draw=white,line width=0.8pt] (3.10,-7.35) -- (2.96,-7.54) -- (2.89,-7.53) -- (2.93,-7.30) -- (2.91,-7.24) -- (2.84,-7.24) -- (2.87,-7.07) -- (2.91,-7.02) -- (3.01,-7.02) -- (3.01,-6.96) -- (3.21,-7.02) -- (3.30,-6.96) -- (3.26,-7.07) -- (3.17,-7.12) -- (3.12,-7.21) -- cycle;
\filldraw[fill=g1,draw=white,line width=0.8pt] (2.82,-6.82) -- (2.75,-6.82) -- (2.68,-6.68) -- (2.70,-6.59) -- (2.63,-6.56) -- (2.56,-6.40) -- (2.57,-6.35) -- (2.68,-6.31) -- (2.68,-6.24) -- (2.71,-6.20) -- (2.86,-6.27) -- (2.87,-6.38) -- (3.00,-6.36) -- (3.07,-6.43) -- (3.01,-6.49) -- (2.91,-6.47) -- (2.93,-6.54) -- (2.89,-6.61) -- (2.84,-6.59) -- (2.82,-6.65) -- (2.84,-6.75) -- (2.91,-6.79) -- (2.98,-6.68) -- (3.01,-6.77) -- (3.07,-6.73) -- (3.08,-6.79) -- (3.21,-6.79) -- (3.30,-6.89) -- (3.30,-6.96) -- (3.21,-7.02) -- (3.16,-6.98) -- (2.87,-6.96) -- (2.87,-6.91) -- (2.80,-6.93) -- cycle;
\filldraw[fill=g1,draw=white,line width=0.8pt] (4.78,-6.03) -- (4.69,-6.08) -- (4.67,-6.06) -- (4.53,-6.08) -- (4.49,-5.99) -- (4.37,-5.90) -- (4.35,-5.96) -- (4.28,-5.98) -- (4.32,-5.85) -- (4.30,-5.76) -- (4.27,-5.75) -- (4.18,-5.78) -- (4.12,-5.75) -- (4.11,-5.64) -- (4.05,-5.59) -- (4.20,-5.46) -- (4.16,-5.31) -- (4.42,-5.15) -- (4.51,-5.01) -- (4.62,-4.94) -- (4.58,-4.78) -- (4.62,-4.57) -- (4.60,-4.51) -- (4.71,-4.35) -- (4.64,-3.98) -- (4.72,-3.91) -- (4.71,-3.95) -- (4.81,-4.12) -- (4.74,-4.21) -- (4.74,-4.27) -- (4.85,-4.27) -- (4.97,-4.41) -- (5.18,-4.46) -- (5.22,-4.37) -- (5.43,-4.39) -- (5.43,-4.46) -- (5.50,-4.51) -- (5.55,-4.62) -- (5.52,-4.71) -- (5.55,-4.76) -- (5.55,-4.85) -- (5.64,-4.88) -- (5.64,-4.95) -- (5.69,-4.94) -- (5.71,-4.97) -- (5.80,-4.94) -- (5.82,-5.04) -- (5.85,-5.02) -- (5.87,-5.08) -- (5.82,-5.39) -- (5.73,-5.43) -- (5.62,-5.31) -- (5.59,-5.32) -- (5.66,-5.41) -- (5.64,-5.62) -- (5.48,-5.73) -- (5.41,-5.85) -- (5.38,-5.80) -- (5.32,-5.85) -- (5.24,-5.82) -- (5.22,-5.92) -- (5.25,-5.96) -- (5.20,-5.99) -- (5.09,-5.94) -- (5.06,-6.05) -- (4.99,-5.99) -- (5.08,-5.82) -- (4.97,-5.73) -- (4.85,-5.90) -- (4.88,-5.94) -- (4.85,-6.08) -- cycle;
\filldraw[fill=g2,draw=white,line width=0.8pt] (3.42,-3.26) -- (3.49,-3.21) -- (3.44,-3.08) -- (3.38,-3.08) -- (3.30,-3.01) -- (3.24,-3.03) -- (3.19,-2.91) -- (3.17,-2.93) -- (3.10,-2.91) -- (3.12,-2.89) -- (3.08,-2.84) -- (3.01,-2.82) -- (3.01,-2.78) -- (3.10,-2.73) -- (3.12,-2.61) -- (3.08,-2.54) -- (3.01,-2.54) -- (3.08,-2.43) -- (3.12,-2.19) -- (3.19,-2.13) -- (3.16,-2.22) -- (3.21,-2.24) -- (3.23,-2.36) -- (3.24,-2.34) -- (3.33,-2.41) -- (3.30,-2.59) -- (3.35,-2.61) -- (3.38,-2.54) -- (3.49,-2.57) -- (3.47,-2.63) -- (3.63,-2.68) -- (3.61,-2.73) -- (3.74,-2.77) -- (3.68,-2.84) -- (3.72,-3.01) -- (3.79,-3.10) -- (3.95,-3.17) -- (4.04,-3.33) -- (3.98,-3.37) -- (4.00,-3.53) -- (4.05,-3.63) -- (4.02,-3.70) -- (3.98,-3.68) -- (3.95,-3.79) -- (3.82,-3.79) -- (3.79,-3.72) -- (3.81,-3.70) -- (3.70,-3.65) -- (3.70,-3.60) -- (3.63,-3.58) -- (3.60,-3.63) -- (3.53,-3.63) -- (3.47,-3.68) -- (3.44,-3.61) -- (3.35,-3.60) -- (3.31,-3.44) -- (3.35,-3.37) -- (3.30,-3.31) -- (3.38,-3.23) -- cycle;
\filldraw[fill=g1,draw=white,line width=0.8pt] (3.21,-8.18) -- (3.08,-8.43) -- (3.10,-8.46) -- (2.93,-8.60) -- (2.87,-8.62) -- (2.78,-8.57) -- (2.75,-8.51) -- (2.71,-8.55) -- (2.73,-8.50) -- (2.64,-8.41) -- (2.68,-8.39) -- (2.66,-8.32) -- (2.71,-8.27) -- (2.70,-8.18) -- (2.77,-8.00) -- (2.75,-7.93) -- (3.10,-7.35) -- (3.12,-7.21) -- (3.17,-7.12) -- (3.26,-7.05) -- (3.35,-6.82) -- (3.31,-6.77) -- (3.35,-6.70) -- (3.33,-6.65) -- (3.38,-6.59) -- (3.53,-6.57) -- (3.61,-6.45) -- (3.79,-6.43) -- (3.86,-6.45) -- (3.90,-6.56) -- (3.86,-6.72) -- (3.90,-6.73) -- (3.86,-6.75) -- (3.86,-6.91) -- (3.79,-6.96) -- (3.82,-7.16) -- (3.74,-7.35) -- (3.82,-7.35) -- (3.91,-7.42) -- (4.00,-7.37) -- (4.05,-7.39) -- (4.11,-7.49) -- (4.05,-7.56) -- (4.05,-7.77) -- (3.86,-7.97) -- (3.88,-8.00) -- (3.82,-8.04) -- (3.84,-8.11) -- (3.72,-8.23) -- (3.58,-8.23) -- (3.60,-8.11) -- (3.51,-8.11) -- (3.47,-8.16) -- (3.44,-8.13) -- (3.40,-8.18) -- (3.35,-8.13) -- (3.24,-8.20) -- cycle;
\filldraw[fill=g1,draw=white,line width=0.8pt] (2.34,-8.27) -- (2.29,-8.39) -- (2.26,-8.41) -- (2.10,-8.39) -- (2.06,-8.28) -- (1.94,-8.32) -- (1.92,-8.39) -- (1.76,-8.34) -- (1.66,-8.41) -- (1.57,-8.34) -- (1.53,-8.37) -- (1.48,-8.27) -- (1.43,-8.32) -- (1.37,-8.30) -- (1.53,-8.09) -- (1.53,-8.13) -- (1.57,-8.13) -- (1.53,-8.07) -- (1.59,-7.98) -- (1.71,-7.93) -- (1.69,-7.77) -- (1.75,-7.77) -- (1.82,-7.72) -- (1.78,-7.69) -- (1.69,-7.76) -- (1.60,-7.76) -- (1.57,-7.65) -- (1.69,-7.56) -- (1.62,-7.54) -- (1.53,-7.67) -- (1.46,-7.58) -- (1.46,-7.49) -- (1.55,-7.44) -- (1.59,-7.46) -- (1.64,-7.40) -- (1.66,-7.47) -- (1.71,-7.44) -- (1.76,-7.51) -- (1.90,-7.49) -- (2.03,-7.60) -- (2.13,-7.63) -- (2.19,-7.60) -- (2.22,-7.63) -- (2.34,-7.46) -- (2.26,-7.24) -- (2.41,-7.17) -- (2.38,-7.14) -- (2.47,-7.07) -- (2.49,-6.98) -- (2.47,-6.94) -- (2.63,-6.86) -- (2.66,-6.77) -- (2.66,-6.68) -- (2.75,-6.82) -- (2.84,-6.84) -- (2.82,-6.93) -- (2.87,-6.91) -- (2.87,-6.96) -- (3.00,-6.98) -- (3.01,-7.02) -- (2.91,-7.02) -- (2.87,-7.07) -- (2.87,-7.21) -- (2.84,-7.23) -- (2.91,-7.24) -- (2.93,-7.30) -- (2.89,-7.53) -- (2.98,-7.56) -- (2.75,-7.93) -- (2.70,-8.30) -- (2.34,-8.20) -- (2.31,-8.23) -- cycle;
\filldraw[fill=g1,draw=white,line width=0.8pt] (7.97,-4.95) -- (8.58,-5.73) -- (8.48,-5.78) -- (8.20,-5.78) -- (8.13,-5.69) -- (8.02,-5.68) -- (8.04,-5.66) -- (7.93,-5.55) -- (7.84,-5.59) -- (7.74,-5.55) -- (7.61,-5.62) -- (7.44,-5.64) -- (7.37,-5.61) -- (7.21,-5.62) -- (7.00,-5.55) -- (6.87,-5.62) -- (6.79,-5.61) -- (6.80,-5.64) -- (6.72,-5.62) -- (6.68,-5.66) -- (6.63,-5.62) -- (6.59,-5.66) -- (6.35,-5.68) -- (6.17,-5.76) -- (6.06,-5.71) -- (6.05,-5.64) -- (5.98,-5.59) -- (6.06,-5.43) -- (6.10,-5.45) -- (6.19,-5.36) -- (6.20,-5.16) -- (6.26,-5.15) -- (6.33,-4.94) -- (6.47,-4.88) -- (6.56,-4.90) -- (6.66,-4.88) -- (6.66,-4.92) -- (7.07,-4.90) -- (7.09,-4.94) -- (7.17,-4.95) -- (7.26,-4.88) -- (7.47,-4.85) -- (7.51,-4.88) -- (7.58,-4.86) -- (7.63,-4.94) -- (7.70,-4.94) -- (7.79,-4.99) -- cycle;
\filldraw[fill=g1,draw=white,line width=0.8pt] (5.50,-6.66) -- (5.43,-6.57) -- (5.50,-6.49) -- (5.52,-6.40) -- (5.53,-6.42) -- (5.59,-6.38) -- (5.68,-6.47) -- (5.78,-6.40) -- (5.92,-6.22) -- (6.03,-6.26) -- (6.06,-6.13) -- (6.10,-6.12) -- (6.03,-6.01) -- (5.89,-5.96) -- (5.92,-5.90) -- (6.01,-5.90) -- (5.98,-5.80) -- (6.01,-5.75) -- (5.98,-5.76) -- (5.92,-5.66) -- (5.99,-5.52) -- (5.98,-5.59) -- (6.05,-5.64) -- (6.05,-5.71) -- (6.17,-5.76) -- (6.35,-5.68) -- (6.59,-5.66) -- (6.63,-5.62) -- (6.68,-5.66) -- (6.72,-5.62) -- (6.80,-5.64) -- (6.79,-5.61) -- (6.87,-5.62) -- (7.00,-5.55) -- (7.21,-5.62) -- (7.37,-5.61) -- (7.44,-5.64) -- (7.61,-5.62) -- (7.76,-5.55) -- (7.84,-5.59) -- (7.93,-5.57) -- (8.04,-5.66) -- (8.02,-5.68) -- (8.09,-5.68) -- (8.18,-5.78) -- (8.23,-5.78) -- (7.95,-6.17) -- (7.84,-6.22) -- (7.76,-6.19) -- (7.47,-6.40) -- (7.46,-6.47) -- (7.31,-6.59) -- (7.24,-6.61) -- (7.14,-6.80) -- (6.54,-7.10) -- (6.47,-7.09) -- (6.31,-7.26) -- (6.08,-7.23) -- (5.94,-7.33) -- (5.92,-7.30) -- (5.89,-7.33) -- (5.82,-7.30) -- (5.78,-7.33) -- (5.64,-7.33) -- (5.39,-7.03) -- (5.34,-6.89) -- (5.39,-6.72) -- (5.43,-6.73) -- cycle;
\filldraw[fill=g2,draw=white,line width=0.8pt] (2.22,-10.98) -- (2.10,-10.91) -- (2.08,-10.95) -- (1.99,-10.96) -- (1.78,-10.91) -- (1.75,-10.65) -- (1.69,-10.61) -- (1.73,-10.43) -- (1.76,-10.43) -- (1.75,-10.38) -- (1.83,-10.28) -- (1.83,-10.24) -- (1.78,-10.21) -- (1.78,-10.12) -- (1.87,-10.08) -- (1.87,-10.03) -- (1.92,-10.03) -- (1.96,-9.94) -- (2.15,-9.87) -- (2.24,-9.94) -- (2.22,-10.01) -- (2.27,-10.10) -- (2.24,-10.21) -- (2.34,-10.31) -- (2.54,-10.31) -- (2.59,-10.26) -- (2.68,-10.26) -- (2.66,-10.22) -- (2.75,-10.26) -- (2.75,-10.21) -- (2.82,-10.24) -- (2.78,-10.29) -- (2.87,-10.38) -- (2.93,-10.36) -- (2.96,-10.42) -- (3.07,-10.38) -- (3.19,-10.49) -- (3.28,-10.51) -- (3.33,-10.49) -- (3.40,-10.52) -- (3.38,-10.58) -- (3.44,-10.61) -- (3.44,-10.73) -- (3.54,-10.72) -- (3.54,-10.75) -- (3.61,-10.75) -- (3.65,-10.80) -- (3.61,-10.82) -- (3.65,-10.89) -- (3.63,-10.93) -- (3.91,-11.02) -- (3.93,-11.09) -- (3.90,-11.09) -- (3.95,-11.10) -- (3.90,-11.14) -- (3.95,-11.18) -- (3.91,-11.23) -- (3.97,-11.21) -- (3.97,-11.25) -- (4.05,-11.28) -- (4.05,-11.25) -- (4.09,-11.30) -- (4.12,-11.23) -- (4.14,-11.26) -- (4.12,-11.30) -- (4.21,-11.26) -- (4.28,-11.30) -- (4.27,-11.35) -- (4.32,-11.39) -- (4.35,-11.35) -- (4.37,-11.40) -- (4.42,-11.39) -- (4.51,-11.47) -- (4.67,-11.53) -- (4.14,-11.67) -- (4.14,-11.62) -- (4.05,-11.58) -- (3.98,-11.49) -- (3.95,-11.51) -- (3.91,-11.46) -- (3.88,-11.47) -- (3.82,-11.35) -- (3.74,-11.39) -- (3.49,-11.21) -- (3.38,-11.28) -- (3.12,-11.19) -- (3.07,-11.12) -- (2.93,-11.12) -- (2.87,-11.05) -- (2.84,-11.05) -- (2.73,-10.91) -- (2.68,-10.88) -- (2.64,-10.89) -- (2.63,-10.84) -- (2.54,-10.84) -- (2.50,-10.79) -- (2.43,-10.84) -- (2.38,-10.82) -- (2.38,-10.96) -- (2.27,-10.95) -- cycle;
\filldraw[fill=g3,draw=white,line width=0.8pt] (7.63,-13.38) -- (7.58,-13.33) -- (7.58,-13.38) -- (7.53,-13.38) -- (7.53,-13.31) -- (7.47,-13.31) -- (7.49,-13.27) -- (7.33,-13.22) -- (7.30,-13.15) -- (7.24,-13.18) -- (7.19,-13.13) -- (7.19,-13.22) -- (7.09,-13.20) -- (7.03,-13.29) -- (6.98,-13.25) -- (6.94,-13.29) -- (6.93,-13.25) -- (6.84,-13.33) -- (6.79,-13.20) -- (6.75,-13.27) -- (6.73,-13.18) -- (6.61,-13.13) -- (6.54,-13.15) -- (6.52,-13.08) -- (6.49,-13.13) -- (6.45,-13.13) -- (6.40,-13.24) -- (6.36,-13.24) -- (6.36,-13.29) -- (6.26,-13.25) -- (6.17,-13.38) -- (6.13,-13.34) -- (6.08,-13.34) -- (5.98,-13.40) -- (5.94,-13.33) -- (5.83,-13.33) -- (5.75,-13.27) -- (5.69,-13.36) -- (5.68,-13.31) -- (5.64,-13.33) -- (5.61,-13.31) -- (5.46,-13.36) -- (5.43,-13.25) -- (5.38,-13.24) -- (5.36,-13.27) -- (5.31,-13.24) -- (5.25,-13.15) -- (5.36,-13.03) -- (5.31,-13.03) -- (5.32,-12.85) -- (5.27,-12.83) -- (5.29,-12.78) -- (5.25,-12.78) -- (5.22,-12.73) -- (5.20,-12.78) -- (5.16,-12.76) -- (5.15,-12.80) -- (5.02,-12.76) -- (4.94,-12.69) -- (4.95,-12.62) -- (5.01,-12.57) -- (4.97,-12.57) -- (4.97,-12.50) -- (4.88,-12.44) -- (4.92,-12.43) -- (4.86,-12.30) -- (4.76,-12.29) -- (4.74,-12.23) -- (4.72,-12.30) -- (4.64,-12.27) -- (4.60,-12.18) -- (4.55,-12.18) -- (4.57,-12.14) -- (4.51,-12.14) -- (4.49,-12.07) -- (4.44,-12.11) -- (4.28,-12.04) -- (4.28,-11.99) -- (4.23,-11.95) -- (4.28,-11.95) -- (4.30,-11.90) -- (4.20,-11.83) -- (4.21,-11.79) -- (4.18,-11.67) -- (4.78,-11.51) -- (4.85,-11.55) -- (4.83,-11.58) -- (4.90,-11.56) -- (4.90,-11.63) -- (4.94,-11.60) -- (4.95,-11.63) -- (5.08,-11.65) -- (5.11,-11.62) -- (5.24,-11.72) -- (5.34,-11.72) -- (5.39,-11.63) -- (5.48,-11.72) -- (5.64,-11.67) -- (5.80,-11.79) -- (5.92,-11.67) -- (6.01,-11.76) -- (6.08,-11.72) -- (6.17,-11.49) -- (6.24,-11.46) -- (6.26,-11.40) -- (6.45,-11.42) -- (6.56,-11.39) -- (6.54,-11.32) -- (6.56,-11.23) -- (6.68,-11.16) -- (6.79,-11.18) -- (7.03,-11.09) -- (7.07,-11.16) -- (7.16,-11.18) -- (7.16,-11.21) -- (7.21,-11.18) -- (7.26,-11.21) -- (7.28,-11.28) -- (7.24,-11.30) -- (7.35,-11.49) -- (7.42,-11.42) -- (7.47,-11.47) -- (7.51,-11.46) -- (7.56,-11.53) -- (7.61,-11.49) -- (7.67,-11.56) -- (7.67,-11.62) -- (7.72,-11.58) -- (7.77,-11.62) -- (7.79,-11.58) -- (7.76,-11.58) -- (7.83,-11.55) -- (7.90,-11.56) -- (7.90,-11.60) -- (7.86,-11.58) -- (7.79,-11.67) -- (7.84,-11.70) -- (7.86,-11.83) -- (7.81,-11.88) -- (7.88,-11.95) -- (7.81,-12.02) -- (7.84,-12.07) -- (7.90,-12.04) -- (7.88,-12.11) -- (7.93,-12.16) -- (8.00,-12.09) -- (7.95,-12.07) -- (7.97,-12.04) -- (8.00,-12.02) -- (8.06,-12.04) -- (8.06,-12.00) -- (8.07,-12.04) -- (8.13,-12.02) -- (8.16,-12.04) -- (8.13,-12.06) -- (8.13,-12.14) -- (8.21,-12.09) -- (8.28,-12.14) -- (8.32,-12.13) -- (8.28,-12.09) -- (8.34,-12.11) -- (8.39,-12.16) -- (8.35,-12.16) -- (8.37,-12.27) -- (8.44,-12.20) -- (8.46,-12.23) -- (8.55,-12.25) -- (8.57,-12.13) -- (8.62,-12.16) -- (8.58,-12.25) -- (8.62,-12.37) -- (8.57,-12.50) -- (8.13,-14.93) -- (8.02,-14.81) -- (7.93,-14.77) -- (7.91,-14.67) -- (7.83,-14.58) -- (7.76,-14.56) -- (7.61,-14.63) -- (7.56,-14.58) -- (7.53,-14.59) -- (7.44,-14.54) -- (8.02,-13.63) -- (8.00,-13.54) -- (7.93,-13.59) -- (7.93,-13.52) -- (7.88,-13.54) -- (7.86,-13.45) -- (7.76,-13.45) -- (7.81,-13.40) -- (7.70,-13.43) -- (7.67,-13.36) -- cycle;
\filldraw[fill=g4,draw=white,line width=0.8pt] (10.68,-9.45) -- (10.82,-9.91) -- (10.88,-9.96) -- (10.84,-10.03) -- (10.88,-10.08) -- (10.66,-10.14) -- (10.61,-9.99) -- (10.68,-9.94) -- (10.65,-9.76) -- (10.61,-9.75) -- (10.63,-9.59) -- (10.54,-9.47) -- (10.47,-9.48) -- (10.42,-9.36) -- (10.45,-9.25) -- (10.38,-9.18) -- (10.35,-9.24) -- (10.24,-9.24) -- (10.21,-9.34) -- (10.15,-9.39) -- (10.05,-9.36) -- (9.98,-9.48) -- (9.92,-9.48) -- (9.94,-9.57) -- (9.85,-9.52) -- (9.84,-9.55) -- (9.68,-9.36) -- (9.61,-9.54) -- (9.66,-9.54) -- (9.64,-9.59) -- (9.71,-9.62) -- (8.62,-9.62) -- (8.50,-9.59) -- (8.41,-9.64) -- (8.32,-9.62) -- (8.25,-9.66) -- (8.04,-9.59) -- (7.91,-9.50) -- (7.98,-9.38) -- (7.97,-9.29) -- (8.07,-9.20) -- (8.04,-9.13) -- (7.90,-9.18) -- (7.81,-9.17) -- (7.81,-9.20) -- (7.72,-9.15) -- (7.63,-9.18) -- (7.61,-9.13) -- (7.51,-9.11) -- (7.40,-8.92) -- (7.26,-8.88) -- (7.24,-8.85) -- (7.76,-8.62) -- (7.76,-8.57) -- (7.83,-8.55) -- (7.81,-8.46) -- (7.84,-8.46) -- (7.79,-8.43) -- (7.83,-8.41) -- (7.86,-8.34) -- (7.90,-8.34) -- (7.88,-8.32) -- (7.93,-8.34) -- (7.91,-8.28) -- (7.97,-8.32) -- (7.93,-8.25) -- (7.98,-8.23) -- (7.93,-8.18) -- (7.97,-8.14) -- (8.13,-8.14) -- (8.11,-8.09) -- (8.20,-8.11) -- (8.18,-8.02) -- (8.23,-8.07) -- (8.21,-8.00) -- (8.25,-8.04) -- (8.27,-8.00) -- (8.34,-8.00) -- (8.30,-7.97) -- (8.34,-8.02) -- (8.37,-7.98) -- (8.44,-7.98) -- (8.44,-7.90) -- (8.53,-7.91) -- (8.57,-7.86) -- (8.60,-7.91) -- (8.69,-7.86) -- (8.71,-7.93) -- (8.73,-7.84) -- (8.76,-7.90) -- (8.78,-7.88) -- (8.81,-7.90) -- (8.80,-7.93) -- (8.87,-7.90) -- (8.92,-7.93) -- (8.97,-7.88) -- (9.02,-7.90) -- (9.02,-7.84) -- (9.02,-7.88) -- (9.08,-7.86) -- (9.11,-7.90) -- (9.10,-7.84) -- (9.15,-7.86) -- (9.18,-7.81) -- (9.27,-7.83) -- (9.25,-7.77) -- (9.36,-7.83) -- (9.34,-7.77) -- (9.39,-7.81) -- (9.39,-7.76) -- (9.47,-7.76) -- (9.47,-7.70) -- (9.50,-7.74) -- (9.47,-7.69) -- (9.55,-7.67) -- (9.55,-7.61) -- (9.68,-7.65) -- (9.68,-7.61) -- (9.73,-7.61) -- (9.73,-7.65) -- (9.78,-7.63) -- (9.78,-7.67) -- (9.82,-7.60) -- (9.84,-7.63) -- (9.91,-7.63) -- (9.92,-7.69) -- (10.12,-7.56) -- (10.17,-7.79) -- (10.21,-7.83) -- (10.29,-7.81) -- (10.38,-8.07) -- (10.43,-8.09) -- (10.47,-8.16) -- (10.40,-8.28) -- (10.36,-8.28) -- (9.98,-8.62) -- (9.98,-8.69) -- (10.03,-8.65) -- (10.19,-8.67) -- (10.24,-8.81) -- (10.31,-8.80) -- (10.56,-9.04) -- (10.58,-9.08) -- (10.54,-9.17) -- (10.59,-9.27) -- (10.65,-9.29) -- (10.63,-9.41) -- cycle;
\filldraw[fill=g1,draw=white,line width=0.8pt] (6.84,-9.62) -- (6.84,-9.75) -- (6.70,-10.03) -- (6.70,-10.10) -- (6.75,-10.19) -- (6.70,-10.14) -- (6.72,-10.17) -- (6.66,-10.17) -- (6.66,-10.24) -- (6.59,-10.31) -- (6.54,-10.29) -- (6.50,-10.36) -- (6.26,-10.59) -- (6.22,-10.58) -- (6.24,-10.56) -- (6.19,-10.51) -- (6.10,-10.54) -- (6.08,-10.47) -- (6.06,-10.49) -- (6.03,-10.45) -- (6.05,-10.40) -- (6.01,-10.42) -- (5.99,-10.35) -- (5.96,-10.35) -- (5.89,-10.21) -- (5.85,-10.22) -- (5.73,-10.12) -- (5.71,-10.15) -- (5.66,-10.12) -- (5.64,-10.17) -- (5.62,-10.12) -- (5.53,-10.14) -- (5.53,-10.26) -- (5.46,-10.26) -- (5.43,-10.35) -- (5.34,-10.36) -- (5.06,-10.15) -- (5.02,-10.01) -- (4.83,-9.76) -- (4.83,-9.13) -- (4.88,-9.04) -- (4.90,-9.08) -- (4.95,-9.04) -- (4.95,-9.10) -- (5.01,-9.08) -- (5.02,-9.04) -- (5.02,-9.10) -- (5.04,-9.08) -- (5.08,-9.10) -- (5.09,-9.06) -- (5.09,-9.10) -- (5.11,-9.06) -- (5.20,-9.08) -- (5.18,-9.04) -- (5.22,-9.02) -- (5.27,-9.02) -- (5.27,-9.06) -- (5.45,-8.95) -- (5.46,-8.97) -- (5.45,-8.94) -- (5.48,-8.88) -- (5.57,-8.85) -- (5.64,-8.88) -- (5.69,-8.83) -- (5.73,-8.88) -- (5.75,-8.83) -- (5.78,-8.87) -- (5.82,-8.76) -- (5.89,-8.78) -- (5.89,-8.74) -- (5.94,-8.74) -- (6.01,-8.67) -- (6.05,-8.69) -- (6.06,-8.65) -- (6.13,-8.71) -- (6.12,-8.62) -- (6.19,-8.60) -- (6.22,-8.65) -- (6.31,-8.67) -- (6.45,-8.64) -- (6.43,-8.62) -- (6.45,-8.60) -- (6.50,-8.65) -- (6.50,-8.58) -- (6.57,-8.64) -- (6.59,-8.62) -- (6.61,-8.65) -- (6.68,-8.64) -- (6.66,-8.62) -- (6.77,-8.64) -- (6.77,-8.58) -- (6.84,-8.64) -- (6.91,-8.57) -- (6.94,-8.62) -- (7.03,-8.58) -- (7.05,-8.62) -- (7.07,-8.58) -- (7.23,-8.64) -- (7.24,-8.69) -- (7.24,-8.62) -- (7.30,-8.69) -- (7.30,-8.64) -- (7.30,-8.69) -- (7.35,-8.67) -- (7.35,-8.71) -- (7.44,-8.67) -- (7.46,-8.73) -- (7.51,-8.67) -- (7.49,-8.62) -- (7.53,-8.69) -- (7.56,-8.64) -- (7.63,-8.69) -- (7.24,-8.85) -- (7.24,-8.88) -- (7.40,-8.92) -- (7.51,-9.11) -- (7.61,-9.13) -- (7.63,-9.18) -- (7.72,-9.15) -- (7.81,-9.20) -- (7.81,-9.17) -- (7.90,-9.18) -- (8.02,-9.13) -- (8.07,-9.18) -- (7.91,-9.36) -- (7.69,-9.39) -- (7.51,-9.48) -- (7.30,-9.45) -- (7.19,-9.55) -- (7.10,-9.57) -- (6.96,-9.68) -- cycle;
\filldraw[fill=g3,draw=white,line width=0.8pt] (8.04,-10.66) -- (8.04,-11.33) -- (8.14,-11.46) -- (8.21,-11.47) -- (8.43,-11.62) -- (8.43,-11.70) -- (8.46,-11.74) -- (8.41,-11.83) -- (8.50,-11.95) -- (8.50,-11.99) -- (8.60,-12.06) -- (8.57,-12.07) -- (8.58,-12.13) -- (8.55,-12.13) -- (8.53,-12.25) -- (8.46,-12.23) -- (8.44,-12.20) -- (8.37,-12.27) -- (8.35,-12.16) -- (8.39,-12.16) -- (8.30,-12.09) -- (8.28,-12.09) -- (8.30,-12.14) -- (8.23,-12.09) -- (8.11,-12.14) -- (8.13,-12.06) -- (8.16,-12.04) -- (8.13,-12.02) -- (8.07,-12.04) -- (8.06,-12.00) -- (8.04,-12.04) -- (7.98,-12.02) -- (7.95,-12.07) -- (8.00,-12.09) -- (7.93,-12.16) -- (7.88,-12.11) -- (7.90,-12.04) -- (7.84,-12.07) -- (7.81,-12.02) -- (7.88,-11.95) -- (7.81,-11.88) -- (7.86,-11.83) -- (7.84,-11.70) -- (7.79,-11.67) -- (7.86,-11.58) -- (7.90,-11.60) -- (7.90,-11.56) -- (7.83,-11.55) -- (7.76,-11.58) -- (7.79,-11.58) -- (7.77,-11.62) -- (7.72,-11.58) -- (7.67,-11.62) -- (7.67,-11.56) -- (7.61,-11.49) -- (7.56,-11.53) -- (7.51,-11.46) -- (7.47,-11.47) -- (7.42,-11.42) -- (7.35,-11.49) -- (7.24,-11.30) -- (7.28,-11.26) -- (7.21,-11.18) -- (7.16,-11.21) -- (7.16,-11.18) -- (7.07,-11.16) -- (7.03,-11.09) -- (6.98,-11.12) -- (6.93,-11.10) -- (6.89,-11.03) -- (6.84,-11.00) -- (6.79,-11.05) -- (6.77,-11.00) -- (6.72,-11.07) -- (6.72,-11.02) -- (6.66,-11.02) -- (6.70,-10.98) -- (6.63,-11.00) -- (6.54,-10.91) -- (6.56,-10.89) -- (6.52,-10.91) -- (6.49,-10.88) -- (6.50,-10.84) -- (6.43,-10.86) -- (6.45,-10.80) -- (6.43,-10.77) -- (6.36,-10.77) -- (6.36,-10.68) -- (6.31,-10.66) -- (6.29,-10.63) -- (6.31,-10.61) -- (6.27,-10.58) -- (6.50,-10.36) -- (6.54,-10.29) -- (6.59,-10.31) -- (6.66,-10.24) -- (6.66,-10.17) -- (6.72,-10.17) -- (6.70,-10.14) -- (6.75,-10.19) -- (6.70,-10.10) -- (6.70,-10.03) -- (6.84,-9.75) -- (6.84,-9.62) -- (6.96,-9.68) -- (7.10,-9.57) -- (7.19,-9.55) -- (7.30,-9.45) -- (7.51,-9.48) -- (7.69,-9.39) -- (7.91,-9.36) -- (7.98,-9.32) -- (7.91,-9.50) -- (8.04,-9.59) -- (8.21,-9.64) -- (8.21,-10.21) -- (8.34,-10.17) -- (8.34,-10.21) -- (8.43,-10.19) -- (8.60,-10.24) -- (8.67,-10.21) -- (8.80,-10.28) -- (8.78,-10.29) -- (8.81,-10.33) -- (8.80,-10.35) -- (8.85,-10.40) -- (8.80,-10.52) -- (8.87,-10.61) -- (8.80,-10.59) -- (8.80,-10.63) -- (8.71,-10.63) -- (8.71,-10.59) -- (8.65,-10.63) -- (8.53,-10.52) -- (8.43,-10.63) -- (8.37,-10.58) -- (8.25,-10.66) -- cycle;
\filldraw[fill=g1,draw=white,line width=0.8pt] (10.19,-6.29) -- (9.99,-6.47) -- (10.01,-6.54) -- (9.98,-6.59) -- (10.03,-6.66) -- (9.98,-7.14) -- (10.05,-7.23) -- (10.03,-7.40) -- (10.12,-7.56) -- (9.92,-7.69) -- (9.91,-7.63) -- (9.84,-7.63) -- (9.82,-7.60) -- (9.78,-7.67) -- (9.78,-7.63) -- (9.73,-7.65) -- (9.73,-7.61) -- (9.68,-7.61) -- (9.68,-7.65) -- (9.55,-7.61) -- (9.55,-7.67) -- (9.47,-7.69) -- (9.50,-7.74) -- (9.47,-7.70) -- (9.47,-7.76) -- (9.39,-7.76) -- (9.39,-7.81) -- (9.34,-7.77) -- (9.36,-7.83) -- (9.25,-7.77) -- (9.27,-7.83) -- (9.18,-7.81) -- (9.15,-7.86) -- (9.10,-7.84) -- (9.11,-7.90) -- (9.08,-7.86) -- (9.02,-7.88) -- (9.02,-7.84) -- (9.02,-7.90) -- (8.97,-7.88) -- (8.92,-7.93) -- (8.87,-7.90) -- (8.80,-7.93) -- (8.81,-7.90) -- (8.78,-7.88) -- (8.76,-7.90) -- (8.73,-7.84) -- (8.71,-7.93) -- (8.69,-7.86) -- (8.60,-7.91) -- (8.57,-7.86) -- (8.53,-7.91) -- (8.44,-7.90) -- (8.44,-7.98) -- (8.37,-7.98) -- (8.34,-8.02) -- (8.34,-7.98) -- (8.30,-7.97) -- (8.34,-8.00) -- (8.27,-8.00) -- (8.25,-8.04) -- (8.21,-8.00) -- (8.23,-8.07) -- (8.18,-8.02) -- (8.20,-8.11) -- (8.11,-8.09) -- (8.13,-8.14) -- (7.97,-8.14) -- (7.93,-8.18) -- (7.98,-8.23) -- (7.93,-8.25) -- (7.97,-8.32) -- (7.91,-8.28) -- (7.93,-8.34) -- (7.88,-8.32) -- (7.90,-8.34) -- (7.86,-8.34) -- (7.83,-8.41) -- (7.79,-8.43) -- (7.84,-8.46) -- (7.81,-8.46) -- (7.83,-8.55) -- (7.76,-8.57) -- (7.76,-8.64) -- (7.63,-8.69) -- (7.56,-8.64) -- (7.53,-8.69) -- (7.49,-8.62) -- (7.51,-8.67) -- (7.46,-8.73) -- (7.44,-8.67) -- (7.35,-8.71) -- (7.35,-8.67) -- (7.30,-8.69) -- (7.31,-8.65) -- (7.30,-8.64) -- (7.30,-8.69) -- (7.24,-8.62) -- (7.24,-8.69) -- (7.23,-8.64) -- (7.14,-8.62) -- (7.12,-6.80) -- (7.23,-6.61) -- (7.31,-6.59) -- (7.46,-6.47) -- (7.47,-6.40) -- (7.76,-6.19) -- (7.84,-6.22) -- (7.97,-6.15) -- (8.23,-5.78) -- (8.48,-5.78) -- (8.65,-5.69) -- (8.69,-5.75) -- (8.76,-5.75) -- (8.92,-5.62) -- (9.29,-5.71) -- (9.45,-5.66) -- (9.59,-5.68) -- (9.73,-5.62) -- (9.84,-5.66) -- (10.03,-5.53) -- (10.26,-5.59) -- (10.33,-5.66) -- (10.31,-5.71) -- (10.38,-5.83) -- (10.21,-5.98) -- (10.17,-6.12) -- (10.21,-6.24) -- cycle;
\node[anchor=south west,font=\scriptsize\bfseries,text=tinta,align=left] at (11.53,-4.50) {Grupo dominante\\(departamentos)};
\fill[g1] (11.53,-5.75) rectangle (12.23,-5.05); \node[anchor=west,font=\scriptsize,text=tinta] at (12.43,-5.40) {G1 · servicio: 21};
\fill[g2] (11.53,-7.00) rectangle (12.23,-6.30); \node[anchor=west,font=\scriptsize,text=tinta] at (12.43,-6.65) {G2 · ataúdes: 7};
\fill[g3] (11.53,-8.25) rectangle (12.23,-7.55); \node[anchor=west,font=\scriptsize,text=tinta] at (12.43,-7.90) {G3 · traslado: 2};
\fill[g4] (11.53,-9.50) rectangle (12.23,-8.80); \node[anchor=west,font=\scriptsize,text=tinta] at (12.43,-9.15) {G4 · compraventa: 2*};
\fill[g7] (11.53,-10.75) rectangle (12.23,-10.05); \node[anchor=west,font=\scriptsize,text=tinta] at (12.43,-10.40) {G7 · serv. públicos: 1};
\end{tikzpicture}

          \caption{Grupo dominante por departamento (V de Cramér 0,34). *Uno es San Andrés}
        \end{figure}
      \end{column}
    \end{columns}

    \begin{columns}[t]
      \begin{column}{0.49\linewidth}
        \begin{figure}
          \centering
          % Generado por gen_tikz.py
\begin{tikzpicture}
\begin{axis}[width=0.78\linewidth,height=8.5cm,scale only axis,xmin=0.6,xmax=12.4,ymin=0,ymax=195,clip=false,
  xtick={1,...,12},xticklabels={E,F,M,A,M,J,J,A,S,O,N,D},ytick={0,50,100,150},yticklabels={0,50,100,150},
  ticklabel style={font=\footnotesize,text=tintamedia},axis x line*=bottom,axis y line*=left,ymajorgrids,grid style={lineafuerte,line width=1pt},
  axis line style={tintamedia,line width=1.2pt},major tick style={tintamedia},
  ylabel={Contratos por mes},ylabel style={font=\scriptsize,text=tinta}]
\fill[naranja,opacity=.16] (axis cs:1.5,0) rectangle (axis cs:4.5,195);
\addplot[gris,line width=1.5pt,opacity=.85] coordinates {(1,56) (2,100) (3,110) (4,138) (5,115) (6,111) (7,89) (8,96) (9,73) (10,74) (11,50) (12,54)};
\addplot[gris,line width=1.5pt,opacity=.85] coordinates {(1,70) (2,147) (3,176) (4,110) (5,90) (6,82) (7,58) (8,50) (9,55) (10,49) (11,48) (12,38)};
\addplot[gris,line width=1.5pt,opacity=.85] coordinates {(1,93) (2,137) (3,157) (4,87) (5,80) (6,52) (7,39) (8,35) (9,46) (10,27) (11,31) (12,25)};
\addplot[gris,line width=1.5pt,opacity=.85] coordinates {(1,64) (2,128) (3,158) (4,80) (5,73) (6,64) (7,51) (8,26) (9,29) (10,22) (11,29) (12,41)};
\addplot[gris,line width=1.5pt,opacity=.85] coordinates {(1,31) (2,107) (3,95) (4,110) (5,81) (6,40) (7,36) (8,33) (9,32) (10,41) (11,37) (12,47)};
\addplot[gris,line width=1.5pt,opacity=.85] coordinates {(1,50) (2,109) (3,100) (4,101) (5,74) (6,49) (7,60) (8,26) (9,27) (10,33) (11,26) (12,34)};
\addplot[noche,line width=3.5pt,mark=*,mark size=4.5pt] coordinates {(1,60.67) (2,121.33) (3,132.67) (4,104.33) (5,85.50) (6,66.33) (7,55.50) (8,44.33) (9,43.67) (10,41.00) (11,36.83) (12,39.83)};
\node[anchor=north,font=\scriptsize\bfseries,text=naranjatexto,align=center,fill=white,fill opacity=.85,text opacity=1,inner sep=2pt] at (axis cs:3,192) {feb--abr: 43\,\%\\de los contratos};
\draw[gris,line width=1.5pt] (axis cs:1.2,212) -- (axis cs:2.0,212); \node[anchor=west,font=\scriptsize,text=tinta] at (axis cs:2.1,212) {cada año, 2020--2025};
\draw[noche,line width=3.5pt] (axis cs:7.6,212) -- (axis cs:8.4,212); \node[anchor=west,font=\scriptsize,text=tinta] at (axis cs:8.5,212) {promedio};
\end{axis}
\end{tikzpicture}

          \caption{Contratos por mes}
        \end{figure}
      \end{column}
      \begin{column}{0.49\linewidth}
        \begin{table}
        \caption{Efectos sobre el valor con igual signo y $p < 0{,}05$ en los tres modelos}
        \vskip0.3em
        {\small
        \begin{tabular}{@{}lccc@{}}
          \toprule
          Frente a & Mediana & Mixto & Gamma \\
          \midrule
          Selección abreviada & $\times$4,97 & $\times$3,28 & $\times$3,51 \\
          Gobernación & $\times$1,94 & $\times$2,50 & $\times$2,39 \\
          SECOP II & $\times$1,61 & $\times$1,27 & $\times$1,26 \\
          \bottomrule
        \end{tabular}\par}
        \vskip0.6em
        {\small\raggedright Frente a mínima cuantía, alcaldía y SECOP I; 3 de 21 efectos. La modalidad depende de la cuantía por ley.\par}
        \end{table}
      \end{column}
    \end{columns}

  \end{block}
  \vfill

  \begin{block}{Conclusiones}
    \begin{itemize}
      \setlength{\itemsep}{0.5em}
      \item Entre 2020 y 2025 el Estado gastó en servicios funerarios unos \textbf{254 mil millones de pesos de 2025} (IC 95\,\%: 225–305), unos 42 mil millones al año.
        \fuente{SECOP I y II (Colombia Compra Eficiente) y muestra etiquetada de 400 contratos; IPC del DANE.}
      \item El contrato promedio sube de 42 a 55 millones \textbf{sin encarecerse}: en SECOP I se mantiene (28 a 27) y en SECOP II baja (203 a 82); sube porque la contratación migró a SECOP II.
        \fuente{SECOP I y II, 4.985 contratos funerarios.}
      \item Las alcaldías firman 9 de cada 10 contratos y contratan de dos formas según la región; de febrero a abril se firma el 43\,\% de los contratos.
        \fuente{SECOP I y II; Marco Geoestadístico Nacional del DANE.}
      \item \textbf{Límites:} $\kappa = 0{,}683 < 0{,}70$, residuos del valor con colas pesadas y grupos que se separan en ejes que el plano no muestra.
    \end{itemize}
  \end{block}
  \vfill

  \begin{block}{Referencias}
    \begin{minipage}[c]{0.50\linewidth}
    \footnotesize
    \begin{itemize}
      \setlength{\itemsep}{0.2em}
      \item Fazekas, M., y Kocsis, G. (2020). Uncovering high-level corruption: Cross-national objective corruption risk indicators using public procurement data. \textit{British Journal of Political Science, 50}(1), 155–164.
      \item Särndal, C.-E., Swensson, B., y Wretman, J. (1992). \textit{Model assisted survey sampling}. Springer.
      \item Husson, F., Lê, S., y Pagès, J. (2011). \textit{Exploratory multivariate analysis by example using R}. Chapman \& Hall/CRC.
    \end{itemize}
    \end{minipage}\hfill
    \begin{minipage}[c]{0.23\linewidth}
      \centering
      \includegraphics[width=\linewidth]{qr_referencias.png}\par\vspace{0.2em}
      {\footnotesize\bfseries\color{tinta}Referencias y Fuentes\par}
    \end{minipage}\hfill
    \begin{minipage}[c]{0.23\linewidth}
      \centering
      \includegraphics[width=\linewidth]{qr_repo.png}\par\vspace{0.2em}
      {\footnotesize\bfseries\color{tinta}Investigación y Código\par}
    \end{minipage}\par\vspace{1pt}
  \end{block}

\end{minipage}
\end{column}

\separatorcolumn
\end{columns}
\end{frame}

\end{document}
